\section{Estudio cohomológico de los morfismos proyectivos}
\label{section:III.2}


\subsection{Cálculos explícitos de ciertos grupos de cohomología}
\label{subsection:III.2.1}

\begin{env}[2.1.1]
\label{III.2.1.1}
\oldpage[III]{95}
Sea $X$ un preesquema y sea $\sh{L}$ un $\sh{O}_X$-módulo invertible; consideremos el anillo graduado \sref[0]{0.5.4.6}
\[
\label{III.2.1.1.1}
  S = \Gamma_*(X,\sh{L}) = \bigoplus_{n\in\bb{Z}} \Gamma(X,\sh{L}^{\otimes n}).
  \tag{2.1.1.1}
\]

Sea $(f_i)_{1\leq i\leq r}$ una familia finita de elementos \emph{homogéneos} de $S$, digamos $f_i\in S_{d_i}$; pongamos $U_i = X_{f_i}$, $U = \bigcup_i U_i$, y denotemos por $\mathfrak{U}$ el recubrimiento $(U_i)$ de $U$.
Para todo $\sh{O}_X$-módulo cuasicoherente $\sh{F}$, ponemos
\[
\label{III.2.1.1.2}
  \HH^\bullet(U,\sh{F}(*)) = \bigoplus_{n\in\bb{Z}} \HH^\bullet(U,\sh{F}\otimes\sh{L}^{\otimes n})
  \tag{2.1.1.2}
\]
\[
\label{III.2.1.1.3}
  \HH^\bullet(\mathfrak{U},\sh{F}(*)) = \bigoplus_{n\in\bb{Z}} \HH^\bullet(\mathfrak{U},\sh{F}\otimes\sh{L}^{\otimes n}).
  \tag{2.1.1.3}
\]

Los grupos abelianos (\hyperref[III.2.1.1.2]{2.1.1.2}) y (\hyperref[III.2.1.1.3]{2.1.1.3}) son \emph{bigraduados}, tomando
\[
  (\HH^\bullet(U,\sh{F}(*)))_{mn} = \HH^m(U,\sh{F}\otimes\sh{L}^{\otimes n})
\]
y una definición análoga para (\hyperref[III.2.1.1.3]{2.1.1.3}).
Para la segunda graduación, está claro que estos grupos son $S$-módulos graduados, como se deduce, por ejemplo, del hecho de que $\sh{F}\mapsto\HH^m(U,\sh{F})$ y $\sh{F}\mapsto\HH^m(\mathfrak{U},\sh{F})$ son funtores.
\end{env}

\begin{env}[2.1.2]
\label{III.2.1.2}
Consideremos ahora el $S$-módulo graduado \sref[0]{0.5.4.6}
\[
\label{III.2.1.2.1}
  M = \Gamma_*(\sh{L},\sh{F}) = \HH^0(X,\sh{F}(*)) = \bigoplus_{n\in\bb{Z}} \Gamma(X,\sh{F}\otimes\sh{L}^{\otimes n}).
  \tag{2.1.2.1}
\]

\oldpage[III]{96}
Si $X$ es un preesquema cuyo espacio subyacente es noetheriano, o un esquema cuasicompacto, se deduce de \sref[I]{I.9.3.1} que, poniendo como de costumbre $U_{i_0 i_1\cdots i_p} = \bigcap_{k=0}^p U_{i_k}$, se tiene, salvo isomorfismo canónico,
\[
  \Gamma(U_{i_0 i_1\cdots i_p},\sh{F}(*)) = \HH^0(U_{i_0 i_1\cdots i_p},\sh{F}(*)) = M_{f_{i_0}f_{i_1}\cdots f_{i_p}}.
\]

Además, con la notación de (\hyperref[III.1.2.2]{1.2.2}), podemos identificar $M_{f_{i_0}f_{i_1}\cdots f_{i_p}}$ con $\varinjlim_n M_{i_0 i_1\cdots i_p}^{(n)}$.
Esta identificación es un isomorfismo de $S$-módulos \emph{graduados} si definimos el grado de un elemento homogéneo $z\in\varinjlim_n M_{i_0 i_1\cdots i_p}^{(n)}$ como sigue: $z$ es la imagen canónica de un elemento homogéneo $x\in M_{i_0 i_1\cdots i_p}^{(n)} = M$ de grado $m$, y tomamos como grado de $z$ el número $m-n(d_{i_0}+d_{i_1}+\cdots+d_{i_p})$.
Teniendo en cuenta la definición de los homomorfismos $\vphi_{kh}:M_{i_0 i_1\cdots i_p}^{(h)}\to M_{i_0 i_1\cdots i_p}^{(k)}$ (\hyperref[III.1.2.2]{1.2.2}), se ve inmediatamente que esta definición no depende del ``representante'' $x$ de $z$ elegido.
Denotando, como en (\hyperref[III.1.2.2]{1.2.2}), por $C_n^p(M)$ el conjunto de aplicaciones alternadas de $[1,r]^{p+1}$ en $M$ (para todo $n$), definimos de manera análoga una estructura de $S$-módulo graduado sobre $\varinjlim_n C_n^p(M)$; como en (\hyperref[III.1.2.2]{1.2.2}), se tiene entonces
\[
\label{III.2.1.2.2}
  C^p(\mathfrak{U},\sh{F}(*)) = \varinjlim_n C_n^p(M)
  \tag{2.1.2.2}
\]
donde el isomorfismo entre ambos miembros \emph{respeta los grados}.
Se tiene también, como en (\hyperref[III.1.2.2]{1.2.2}),
\[
\label{III.2.1.2.3}
  C^p(\mathfrak{U},\sh{F}(*)) = C^{p+1}((\mathbf{f}),M) = \varinjlim_n K^{p+1}(\mathbf{f}^n,M)
  \tag{2.1.2.3}
\]
donde el isomorfismo \emph{conserva los grados}: el grado de un elemento de $\varinjlim_n K^{p+1}(\mathbf{f}^n,M)$ que sea la imagen canónica de una cocadena $g\in K^{p+1}(\mathbf{f}^n,M)$ cuyos valores $g(i_0,\dots,i_p)$ pertenezcan a una misma componente homogénea $M_m$ de $M$ es $m-n(d_{i_0}+\cdots+d_{i_p})$, independientemente de la elección de esta cocadena como representante del elemento considerado.

Como los isomorfismos anteriores son compatibles con los operadores de coborde, se concluye, como en (\hyperref[III.1.2.2]{1.2.2}), que:
\end{env}

\begin{proposition}[2.1.3]
\label{III.2.1.3}
Sea $X$ un preesquema cuyo espacio subyacente es noetheriano, o un esquema cuasicompacto.
Existe un isomorfismo canónico, funtorial en $\sh{F}$,
\[
\label{III.2.1.3.1}
  \HH^p(\mathfrak{U},\sh{F}(*)) \simto \HH^{p+1}((\mathbf{f}),M)
  \qquad\text{para todo } p\geq 1.
  \tag{2.1.3.1}
\]
Además, existe una sucesión exacta, funtorial en $\sh{F}$,
\[
\label{III.2.1.3.2}
  0 \to \HH^0((\mathbf{f}),M) \to M \to \HH^0(\mathfrak{U},\sh{F}(*)) \to \HH^1((\mathbf{f}),M) \to 0.
  \tag{2.1.3.2}
\]
Todos los homomorfismos introducidos son de grado $0$ para las estructuras de $S$-módulo graduado, donde $S$ es el anillo (\hyperref[III.2.1.1.1]{2.1.1.1}).
\end{proposition}

\begin{corollary}[2.1.4]
\label{III.2.1.4}
\oldpage[III]{97}
Si $X$ es un esquema cuasicompacto y los $U_i=X_{f_i}$ son afines, existe un isomorfismo canónico, funtorial en $\sh{F}$ y de grado $0$,
\[
\label{III.2.1.4.1}
  \HH^p(U,\sh{F}(*)) \simto \HH^{p+1}((\mathbf{f}),M)
  \qquad\text{para todo } p\geq 1,
  \tag{2.1.4.1}
\]
y una sucesión exacta, funtorial en $\sh{F}$,
\[
\label{III.2.1.4.2}
  0 \to \HH^0((\mathbf{f}),M) \to M \to \HH^0(U,\sh{F}(*)) \to \HH^1((\mathbf{f}),M) \to 0,
  \tag{2.1.4.2}
\]
en la que todos los homomorfismos son de grado $0$.
\end{corollary}

\begin{proof}
Basta aplicar \sref{III.1.4.1} al resultado de \sref{III.2.1.3}.
\end{proof}

La proposición ``local'' análoga a \sref{III.2.1.3} es la siguiente.

\begin{proposition}[2.1.5]
\label{III.2.1.5}
Sea $S$ un anillo graduado positivamente, sean los $f_i$ ($1\leq i\leq r$) elementos homogéneos de $S_+$ de grados $d_i$, y sea $M$ un $S$-módulo graduado.
Sea $X=\Proj(S)$ el espectro primo homogéneo de $S$, y pongamos $U_i=\mathrm{D}_+(f_i)$, $U=\bigcup_i U_i$ y $\HH^\bullet(U,\widetilde{M}(*))=\bigoplus_{n\in\bb{Z}}\HH^\bullet(U,(M(n))^\sim)$.
Existen isomorfismos canónicos, funtoriales en $M$ y de grado $0$ para las estructuras de $S$-módulo graduado,
\[
\label{III.2.1.5.1}
  \HH^p(U,\widetilde{M}(*)) \simto \HH^{p+1}((\mathbf{f}),M)
  \qquad\text{para todo } p\geq 1,
  \tag{2.1.5.1}
\]
y una sucesión exacta, funtorial en $M$,
\[
\label{III.2.1.5.2}
  0 \to \HH^0((\mathbf{f}),M) \to M \to \HH^0(U,\widetilde{M}(*)) \to \HH^1((\mathbf{f}),M) \to 0,
  \tag{2.1.5.2}
\]
en la que todos los homomorfismos son de grado $0$.
\end{proposition}

\begin{proof}
En efecto, por definición \sref[II]{II.2.5.2},
\[
  \Gamma(U_{i_0i_1\cdots i_p},(M(n))^\sim)
  =(M_{f_{i_0}\cdots f_{i_p}})_n,
\]
y, por tanto, $\Gamma(U_{i_0i_1\cdots i_p},\widetilde{M}(*))=M_{f_{i_0}\cdots f_{i_p}}$.
El resto del razonamiento es el mismo que en la demostración de \sref{III.2.1.4}, teniendo en cuenta que $X$ es un esquema.
\end{proof}

\begin{env}[2.1.6]
\label{III.2.1.6}
\emph{Observaciones.} ---
(i) Bajo las hipótesis de \sref{III.2.1.5}, los funtores $\Gamma(U_{i_0i_1\cdots i_p},\widetilde{M}(*))$ son exactos en $M$, por \sref[0]{0.1.3.2}.
El mismo razonamiento que en (\hyperref[III.1.2.5]{1.2.5}) muestra, por tanto, que, si $0\to M'\to M\to M''\to0$ es una sucesión exacta de $S$-módulos graduados cuyos homomorfismos tienen grado $0$, entonces para todo $p\geq0$ existen diagramas conmutativos
\[
\label{III.2.1.6.1}
  \xymatrix{
    \HH^p(U,\widetilde{M''}(*)) \ar[r]^-\partial \ar[d] &
    \HH^{p+1}(U,\widetilde{M'}(*)) \ar[d] \\
    \HH^{p+1}((\mathbf{f}),M'') \ar[r]^-\partial &
    \HH^{p+2}((\mathbf{f}),M')
  }
  \tag{2.1.6.1}
\]

(ii) La Proposición~\sref{III.2.1.5} es especialmente útil cuando $S$ es una $A$-álgebra engendrada por un número finito de elementos de grado $1$, con $A$ noetheriano; en ese caso, todo $\sh{O}_X$-módulo cuasicoherente es de la forma $\widetilde{M}$ \sref[II]{II.2.7.7}.
\end{env}

\begin{env}[2.1.7]
\label{III.2.1.7}
\oldpage[III]{98}
Aplicamos ahora \sref{III.2.1.5} al caso $S=A[\mathrm{T}_0,\dots,\mathrm{T}_r]$, donde $A$ es un anillo arbitrario y las $\mathrm{T}_i$ son indeterminadas, tomando $M=S$ y $f_i=\mathrm{T}_i$.
Quedamos así reducidos esencialmente a calcular $\HH^\bullet((\mathbf{T}),S)$, donde $\mathbf{T}=(\mathrm{T}_i)_{0\leq i\leq r}$.
\end{env}

\begin{lemma}[2.1.8]
\label{III.2.1.8}
Si $S=A[\mathrm{T}_0,\dots,\mathrm{T}_r]$, entonces, con $\mathbf{T}=(\mathrm{T}_i)_{0\leq i\leq r}$,
\[
\label{III.2.1.8.1}
  \HH^i(\mathbf{T}^n,S)=0
  \qquad\text{si }i\neq r+1,
  \tag{2.1.8.1}
\]
\[
\label{III.2.1.8.2}
  \HH^{r+1}(\mathbf{T}^n,S)=S/(\mathbf{T}^n).
  \tag{2.1.8.2}
\]
Por tanto, el $A$-módulo $\HH^{r+1}(\mathbf{T}^n,S)$ tiene una $A$-base formada por las clases módulo $(\mathbf{T}^n)$ de los monomios
\[
  \mathrm{T}_{\mathbf{p}}=\mathrm{T}_0^{p_0}\cdots\mathrm{T}_r^{p_r},
  \qquad \mathbf{p}=(p_0,\dots,p_r),\quad 0\leq p_i<n\ \text{para todo }i.
\]
\end{lemma}

\begin{proof}
Esto se deduce inmediatamente de (\hyperref[III.1.1.3.5]{1.1.3.5}) y de la Proposición~\sref{III.1.1.4}, cuyas hipótesis se satisfacen trivialmente.
\end{proof}

\begin{env}[2.1.9]
\label{III.2.1.9}
Pasando al límite inductivo respecto de $n$, las relaciones (\hyperref[III.2.1.8.1]{2.1.8.1}) dan $\HH^i((\mathbf{T}),S)=0$ para $i\neq r+1$.
Para $i=r+1$, el sistema inductivo consta de los módulos $S/(\mathbf{T}^n)$, siendo el homomorfismo
\[
  \vphi_{nm}:S/(\mathbf{T}^n)\longrightarrow S/(\mathbf{T}^m)
  \qquad (0\leq n\leq m)
\]
la multiplicación por $(\mathrm{T}_0\cdots\mathrm{T}_r)^{m-n}$.
Para $n\geq\sup(p_i)_{0\leq i\leq r}$, denotemos por $\xi_{\mathbf{p}}^{(n)}=\xi_{p_0,\dots,p_r}^{(n)}$ la clase de $\mathrm{T}_0^{n-p_0}\cdots\mathrm{T}_r^{n-p_r}$ módulo $(\mathbf{T}^n)$.
Entonces $\vphi_{nm}(\xi_{\mathbf{p}}^{(n)})=\xi_{\mathbf{p}}^{(m)}$, de modo que estos elementos tienen la misma imagen canónica $\xi_{\mathbf{p}}=\xi_{p_0,\dots,p_r}$ en el límite inductivo $\HH^{r+1}((\mathbf{T}),S)$.
Por la definición del grado en (\hyperref[III.2.1.2]{2.1.2}), el grado de $\xi_{\mathbf{p}}$ es $-|\mathbf{p}|=-(p_0+p_1+\cdots+p_r)$.
Está claro que los $\xi_{\mathbf{p}}^{(n)}$ con $0<p_i\leq n$ para $0\leq i\leq r$ forman una base de $S/(\mathbf{T}^n)$.
Se deduce, pues, inmediatamente de \sref{III.2.1.8}:
\end{env}

\begin{corollary}[2.1.10]
\label{III.2.1.10}
Con la notación de \sref{III.2.1.8},
\[
\label{III.2.1.10.1}
  \HH^i((\mathbf{T}),S)=0
  \qquad\text{para }i\neq r+1,
  \tag{2.1.10.1}
\]
y $\HH^{r+1}((\mathbf{T}),S)$ es un $A$-módulo libre con base los elementos $\xi_{p_0,\dots,p_r}$ tales que $p_i>0$ para $0\leq i\leq r$.
\end{corollary}

\begin{env}[2.1.11]
\label{III.2.1.11}
\emph{Observación.} ---
Sea $N$ un $A$-módulo cualquiera y sea $M=S\otimes_A N$.
El razonamiento de \sref{III.2.1.8} muestra más generalmente que
\[
\label{III.2.1.11.1}
  \HH^i(\mathbf{T}^n,M)=0
  \qquad\text{si }i\neq r+1,
  \tag{2.1.11.1}
\]
\[
\label{III.2.1.11.2}
  \HH^{r+1}(\mathbf{T}^n,M)=(S/(\mathbf{T}^n))\otimes_A N.
  \tag{2.1.11.2}
\]
En efecto, esta última fórmula se deduce directamente de (\hyperref[III.1.1.3.5]{1.1.3.5}).
Además, está claro que
\[
  M/(\mathrm{T}_0^nM+\cdots+\mathrm{T}_{i-1}^nM)
  =\bigl(S/(\mathrm{T}_0^nS+\cdots+\mathrm{T}_{i-1}^nS)\bigr)\otimes_A N,
\]
puesto que el ideal $\mathrm{T}_0^nS+\cdots+\mathrm{T}_{i-1}^nS$ es un sumando directo del $A$-módulo $S$.
Esto permite aplicar \sref{III.1.1.4} a $M$ y da (\hyperref[III.2.1.11.1]{2.1.11.1}).
\end{env}

Combinando \sref{III.2.1.10} y \sref{III.2.1.5}, se obtiene:

\begin{proposition}[2.1.12]
\label{III.2.1.12}
Sea $A$ un anillo, sea $r$ un entero $>0$ y sea $X=\bb{P}_A^r$ \sref[II]{II.4.1.1}.
Entonces:
\begin{enumerate}
  \item[{\rm(i)}] $\HH^i(X,\sh{O}_X(*))=0$ para $i\neq0,r$.
  \item[{\rm(ii)}] El homomorfismo canónico $\alpha:S\to\HH^0(X,\sh{O}_X(*))$ \sref[II]{II.2.6.2} es biyectivo.
\oldpage[III]{99}
  \item[{\rm(iii)}] $\HH^r(X,\sh{O}_X(*))$ es un $A$-módulo libre con base los elementos $\xi_{p_0,\dots,p_r}$, donde $p_i>0$ para $0\leq i\leq r$; el elemento $\xi_{p_0,\dots,p_r}$ tiene grado $-|\mathbf{p}|=-(p_0+\cdots+p_r)$, y
  \[
    \mathrm{T}_i\xi_{p_0,\dots,p_r}
    =\xi_{p_0,\dots,p_{i-1},p_i-1,p_{i+1},\dots,p_r}.
  \]
\end{enumerate}
\end{proposition}

\begin{proof}
En la sucesión exacta (\hyperref[III.2.1.5.2]{2.1.5.2}) aplicada a $M=S=A[\mathrm{T}_0,\dots,\mathrm{T}_r]$, se tiene $\HH^0((\mathbf{T}),S)=0$ y $\HH^1((\mathbf{T}),S)=0$ por \sref{III.2.1.10.1}; además, la Proposición~\sref{III.2.1.5} se aplica con $U=X$, puesto que $X$ es la unión de los $\mathrm{D}_+(\mathrm{T}_i)$ \sref[II]{II.2.3.14}.
Queda identificar la aplicación $S\to\HH^0(X,\sh{O}_X(*))$ de (\hyperref[III.2.1.5.2]{2.1.5.2}) con la aplicación canónica $\alpha$.
Esto se deduce de la identificación canónica de $\HH^0(U,\sh{O}_X(*))$ y $\HH^0(\mathfrak{U},\sh{O}_X(*))$.
\end{proof}

\begin{corollary}[2.1.13]
\label{III.2.1.13}
Los únicos pares $(i,n)$ para los cuales puede ocurrir que $\HH^i(X,\sh{O}_X(n))\neq0$ son $i=0$, $n\geq0$, e $i=r$, $n\leq-(r+1)$.
\end{corollary}

Si $A\neq0$, entonces se tiene efectivamente $\HH^i(X,\sh{O}_X(n))\neq0$ para todos los pares enumerados en \sref{III.2.1.13}.
Esto se deduce de \sref{III.2.1.12}, puesto que $S_n\neq0$ para todo grado $n\geq0$.

En las aplicaciones de este capítulo, utilizaremos principalmente el siguiente resultado menos preciso.

\begin{corollary}[2.1.14]
\label{III.2.1.14}
Los $A$-módulos $\HH^i(X,\sh{O}_X(n))$ son libres de tipo finito; si $i>0$, se anulan para $n>0$.
\end{corollary}

\begin{proposition}[2.1.15]
\label{III.2.1.15}
Sea $Y$ un preesquema, sea $\sh{E}$ un $\sh{O}_Y$-módulo localmente libre de rango $r+1$, sea $X=\bb{P}(\sh{E})$ el fibrado proyectivo definido por $\sh{E}$, y sea $f:X\to Y$ el morfismo estructural.
Los únicos valores de $i$ y $n$ para los cuales puede ocurrir que $\RR^if_*(\sh{O}_X(n))\neq0$ son $i=0$, $n\geq0$, e $i=r$, $n\leq-(r+1)$.
Además, el homomorfismo canónico \sref[II]{II.3.3.2}
\[
  \alpha:\bb{S}_{\sh{O}_Y}(\sh{E})
  \longrightarrow \Gamma_*(\sh{O}_X)
  =\RR^0f_*(\sh{O}_X(*))
  =\bigoplus_{n\in\bb{Z}}f_*(\sh{O}_X(n))
\]
es un isomorfismo.
\end{proposition}

\begin{proof}
La cuestión es local sobre $Y$, por lo que podemos suponer que $Y$ es afín de anillo $A$ y que $\sh{E}=\widetilde{E}$, donde $E=A^{r+1}$.
Quedamos entonces reducidos inmediatamente a \sref{III.2.1.12}, teniendo en cuenta \sref{III.1.4.11}.
\end{proof}

\begin{env}[2.1.16]
\label{III.2.1.16}
\emph{Observación.} ---
Más adelante completaremos los resultados de \sref{III.2.1.15} demostrando las proposiciones siguientes.
Pongamos
\[
  \boldsymbol{\omega}=f^*\bigl(\bigwedge^{r+1}\sh{E}\bigr)(-r-1),
\]
que es un $\sh{O}_X$-módulo invertible.
Entonces:
\begin{enumerate}
  \item[{\rm(i)}] Existe un isomorfismo canónico
  \[
  \label{III.2.1.16.1}
    \rho:\sh{O}_Y \simto \RR^rf_*(\boldsymbol{\omega}).
    \tag{2.1.16.1}
  \]
  \item[{\rm(ii)}] El apareamiento dado por el producto cup \sref[0]{0.12.2.2}
  \[
  \label{III.2.1.16.2}
    \RR^rf_*(\sh{O}_X(n))
    \times \RR^0f_*(\boldsymbol{\omega}(-n))
    \longrightarrow \RR^rf_*(\boldsymbol{\omega}),
    \tag{2.1.16.2}
  \]
  compuesto con $\rho^{-1}$, define un \emph{isomorfismo} de $\RR^rf_*(\sh{O}_X(n))$ sobre el \emph{dual} del $\sh{O}_Y$-módulo localmente libre
  \[
    \RR^0f_*(\boldsymbol{\omega}(-n))
    =\bigl(\bigwedge^{r+1}\sh{E}\bigr)
      \otimes_{\sh{O}_Y}
      \bigl(\bb{S}_{\sh{O}_Y}(\sh{E})\bigr)_{-n}.
  \]
\end{enumerate}
\end{env}

\subsection{El teorema fundamental de los morfismos proyectivos}
\label{subsection:III.2.2}
\oldpage[III]{100}

\begin{theorem}[2.2.1]
\label{III.2.2.1}
(Serre).
Sea $Y$ un preesquema noetheriano, sea $f:X\to Y$ un morfismo propio y sea $\sh{L}$ un $\sh{O}_X$-módulo invertible amplio respecto de $f$.
Para todo $\sh{O}_X$-módulo $\sh{F}$, pongamos
\[
  \sh{F}(n)=\sh{F}\otimes_{\sh{O}_X}\sh{L}^{\otimes n}
  \qquad(n\in\bb{Z}).
\]
Entonces, para todo $\sh{O}_X$-módulo coherente $\sh{F}$:
\begin{enumerate}
  \item[{\rm(i)}] los $\sh{O}_Y$-módulos $\RR^qf_*(\sh{F})$ son coherentes;
  \item[{\rm(ii)}] existe un entero $N$ tal que, para $n\geq N$,
  \[
    \RR^qf_*(\sh{F}(n))=0
    \qquad\text{para todo }q>0;
  \]
  \item[{\rm(iii)}] existe un entero $N$ tal que, para $n\geq N$, el homomorfismo canónico
  \[
    f^*\bigl(f_*(\sh{F}(n))\bigr)\longrightarrow\sh{F}(n)
  \]
  es sobreyectivo.
\end{enumerate}
\end{theorem}

\begin{proof}
Observemos primero que, si el teorema es cierto tras sustituir $\sh{L}$ por $\sh{L}^{\otimes d}$, donde $d>0$, entonces es cierto en su forma original.
En efecto, escribiendo
\[
  \sh{F}(n)
  =\bigl(\sh{F}\otimes\sh{L}^{\otimes r}\bigr)
   \otimes\sh{L}^{\otimes hd},
   \qquad h>0,\quad 0\leq r<d,
\]
la hipótesis da, para cada $r$, un entero $N_r$ tal que (ii) y (iii) se cumplen para $\sh{F}\otimes\sh{L}^{\otimes r}$ siempre que $h\geq N_r$.
Tomando para $N$ el mayor de los enteros $dN_r$, las afirmaciones (ii) y (iii) se cumplen, por tanto, para $n\geq N$.

Podemos, por consiguiente, suponer que $\sh{L}$ es muy amplio respecto de $f$ \sref[II]{II.4.6.11}.
Existe entonces una inmersión abierta dominante sobre $Y$
\[
  i:X\longrightarrow P,
  \qquad P=\Proj(\sh{S}),
\]
donde $\sh{S}$ es una $\sh{O}_Y$-álgebra cuasicoherente graduada en grados positivos, $\sh{S}_1$ es de tipo finito y engendra $\sh{S}$, y $\sh{L}$ es isomorfo a $i^*(\sh{O}_P(1))$ \sref[II]{II.4.4.7}.
Como $f$ es propio, $i$ también es propio \sref[II]{II.5.4.4} y, por tanto, es un isomorfismo $X\simto P$.
Quedamos así reducidos al caso $X=\Proj(\sh{S})$ y $\sh{L}=\sh{O}_X(1)$.
El teorema es entonces consecuencia de la proposición siguiente.
\end{proof}

\begin{proposition}[2.2.2]
\label{III.2.2.2}
Sea $A$ un anillo noetheriano y sea $S$ una $A$-álgebra graduada en grados positivos tal que el $A$-módulo $S_1$ posee un sistema de $r+1$ generadores y engendra el álgebra $S$.
Pongamos $X=\Proj(S)$.
Para todo $\sh{O}_X$-módulo coherente $\sh{F}$:
\begin{enumerate}
  \item[{\rm(i)}] los $A$-módulos $\HH^q(X,\sh{F})$ son de tipo finito;
  \item[{\rm(ii)}] $\HH^q(X,\sh{F})=0$ para $q>r$;
  \item[{\rm(iii)}] existe un entero $N$ tal que, para $n\geq N$,
  \[
    \HH^q(X,\sh{F}(n))=0
    \qquad\text{para todo }q>0;
  \]
  \item[{\rm(iv)}] existe un entero $N$ tal que, para $n\geq N$, $\sh{F}(n)$ está engendrado por sus secciones sobre $X$.
\end{enumerate}
\end{proposition}

\begin{proof}
Mostremos primero cómo \sref{III.2.2.2} implica \sref{III.2.2.1}.
En el caso particular $X=\Proj(\sh{S})$ al que se redujo anteriormente \sref{III.2.2.1}, $Y$ es cuasicompacto.
Puede, por tanto, recubrirse por un número finito de abiertos afines $U_\alpha$, de anillos noetherianos, tales que la restricción de $\sh{S}_1$ a cada $U_\alpha$ esté engendrada por un número finito de secciones sobre $U_\alpha$.
Una vez conocida \sref{III.2.2.2}, basta tomar en las partes (ii) y (iii) de \sref{III.2.2.1} el mayor de los enteros correspondientes a los $U_\alpha$, teniendo en cuenta \sref{III.1.4.11} y \sref[II]{II.3.4.7}.

Para demostrar \sref{III.2.2.2}, observemos que $X$ se identifica con un subpreesquema cerrado de $P=\bb{P}_A^r$ \sref[II]{II.3.6.2}.
Además, si $j:X\to P$ es la inyección canónica, entonces $j_*(\sh{F})$ es un $\sh{O}_P$-módulo coherente y
\[
  j_*(\sh{F}(n))=(j_*\sh{F})(n)
\]
por \sref[II]{II.3.4.5} y \sref[II]{II.3.5.2}.
En virtud de (G, II, Corolario del Teorema~4.9.1), quedamos, por tanto, reducidos a demostrar \sref{III.2.2.2} en el caso particular
\[
  X=\bb{P}_A^r,
  \qquad S=A[T_0,\ldots,T_r].
\]

\oldpage[III]{101}

Como $X$ está recubierto por los $r+1$ abiertos afines $D_+(T_i)$, la afirmación (ii) se deduce de \sref{III.1.4.12}.
La afirmación (iv) ya fue demostrada en \sref[II]{II.2.7.9}.

Demostremos ahora simultáneamente (i) y (iii).
Se cumplen para $\sh{F}=\sh{O}_X(m)$ por \sref{III.2.1.13} y, por tanto, también cuando $\sh{F}$ es suma directa de un número finito de $\sh{O}_X$-módulos de la forma $\sh{O}_X(m_j)$.
Son trivialmente ciertas para $q>r$ por (ii).
Procedemos por inducción descendente sobre $q$.

El haz $\sh{F}$ es isomorfo a un cociente de una suma directa $\sh{E}$ de un número finito de haces $\sh{O}_X(m_j)$ \sref[II]{II.2.7.10}.
Existe, pues, una sucesión exacta
\[
  0\longrightarrow\sh{R}\longrightarrow\sh{E}
   \longrightarrow\sh{F}\longrightarrow0,
\]
donde $\sh{R}$ es coherente \sref[0]{0.5.3.3} y $\sh{E}$ satisface (i) y (iii).
Como el cambio de grado es exacto, para todo $n\in\bb{Z}$ se tiene también una sucesión exacta
\[
  0\longrightarrow\sh{R}(n)\longrightarrow\sh{E}(n)
   \longrightarrow\sh{F}(n)\longrightarrow0,
\]
y, por tanto, una sucesión exacta de cohomología
\[
  \HH^{q-1}(X,\sh{E}(n))
  \longrightarrow\HH^{q-1}(X,\sh{F}(n))
  \longrightarrow\HH^q(X,\sh{R}(n)).
\]
Como $\sh{E}(n)$ es una suma directa de haces $\sh{O}_X(n+m_j)$ \sref[II]{II.2.5.14}, el módulo $\HH^{q-1}(X,\sh{E}(n))$ es de tipo finito.
Análogamente, el módulo $\HH^q(X,\sh{R}(n))$ es de tipo finito por la hipótesis de inducción.
Como $A$ es noetheriano, se deduce que $\HH^{q-1}(X,\sh{F}(n))$ es de tipo finito, lo que demuestra (i).

Por la hipótesis de inducción podemos elegir $N$ de modo que $\HH^q(X,\sh{R}(n))=0$ para $n\geq N$.
Podemos elegir asimismo $N$ de modo que $\HH^{q-1}(X,\sh{E}(n))=0$ para $n\geq N$, puesto que $\sh{E}$ satisface (iii).
La sucesión exacta de cohomología da entonces
\[
  \HH^{q-1}(X,\sh{F}(n))=0
  \qquad(n\geq N),
\]
lo que termina la demostración.
\end{proof}

\begin{corollary}[2.2.3]
\label{III.2.2.3}
Bajo las hipótesis de \sref{III.2.2.1}, sea
\[
  \sh{F}\longrightarrow\sh{G}\longrightarrow\sh{H}
\]
una sucesión exacta de $\sh{O}_X$-módulos coherentes.
Existe un entero $N$ tal que, para $n\geq N$, la sucesión
\[
  f_*(\sh{F}(n))
  \longrightarrow f_*(\sh{G}(n))
  \longrightarrow f_*(\sh{H}(n))
\]
es exacta.
\end{corollary}

\begin{proof}
Sean $\sh{F}'$, $\sh{G}'$ y $\sh{G}''$, respectivamente, el núcleo, la imagen y el conúcleo de $\sh{F}\to\sh{G}$.
Entonces $\sh{G}'$ es el núcleo y $\sh{G}''$ la imagen de $\sh{G}\to\sh{H}$; sea $\sh{H}''$ el conúcleo de este último homomorfismo.
Todos estos $\sh{O}_X$-módulos son coherentes \sref[0]{0.5.3.4}.
Como el cambio de grado es exacto, basta demostrar que, para $n$ suficientemente grande, cada una de las sucesiones
\begin{align*}
  0&\longrightarrow f_*(\sh{F}'(n))
    \longrightarrow f_*(\sh{F}(n))
    \longrightarrow f_*(\sh{G}'(n))
    \longrightarrow0,\\
  0&\longrightarrow f_*(\sh{G}'(n))
    \longrightarrow f_*(\sh{G}(n))
    \longrightarrow f_*(\sh{G}''(n))
    \longrightarrow0,\\
  0&\longrightarrow f_*(\sh{G}''(n))
    \longrightarrow f_*(\sh{H}(n))
    \longrightarrow f_*(\sh{H}''(n))
    \longrightarrow0
\end{align*}
es exacta.
Por consiguiente, podemos suponer que
\[
  0\longrightarrow\sh{F}\longrightarrow\sh{G}
   \longrightarrow\sh{H}\longrightarrow0
\]
es exacta.
La sucesión de cohomología es entonces
\[
  0\longrightarrow f_*(\sh{F}(n))
  \longrightarrow f_*(\sh{G}(n))
  \longrightarrow f_*(\sh{H}(n))
  \longrightarrow \RR^1f_*(\sh{F}(n))
  \longrightarrow\cdots,
\]
y la conclusión se deduce de \sref{III.2.2.1}, (ii).
\end{proof}

\oldpage[III]{102}

\begin{corollary}[2.2.4]
\label{III.2.2.4}
Sea $Y$ un preesquema noetheriano, sea $f:X\to Y$ un morfismo de tipo finito y sea $\sh{L}$ un $\sh{O}_X$-módulo invertible amplio respecto de $f$.
Para todo $\sh{O}_X$-módulo $\sh{F}$, pongamos
\[
  \sh{F}(n)=\sh{F}\otimes_{\sh{O}_X}\sh{L}^{\otimes n}
  \qquad(n\in\bb{Z}).
\]
Sea
\[
  \sh{F}\longrightarrow\sh{G}\longrightarrow\sh{H}
\]
una sucesión exacta de $\sh{O}_X$-módulos coherentes tal que el soporte de $\operatorname{Im}(\sh{F}\to\sh{G})$ sea propio sobre $Y$ \sref[II]{II.5.4.10}.
Entonces existe un entero $N$ tal que, para $n\geq N$, la sucesión
\[
  f_*(\sh{F}(n))
  \longrightarrow f_*(\sh{G}(n))
  \longrightarrow f_*(\sh{H}(n))
\]
es exacta.
\end{corollary}

\begin{proof}
El razonamiento del comienzo de \sref{III.2.2.1} muestra que, si el corolario es cierto para $\sh{L}^{\otimes d}$, donde $d>0$, entonces es cierto para $\sh{L}$.
Podemos, por tanto, suponer que $\sh{L}$ es muy amplio respecto de $f$ \sref[II]{II.4.6.11}.
Podemos entonces identificar $X$ con un subconjunto abierto de un $Y$-preesquema
\[
  Z=\Proj(\sh{S}),
\]
donde $\sh{S}$ es una $\sh{O}_Y$-álgebra cuasicoherente graduada en grados positivos, $\sh{S}_1$ es de tipo finito y engendra $\sh{S}$, y
\[
  \sh{L}=i^*(\sh{O}_Z(1)),
\]
siendo $i:X\to Z$ la inmersión canónica \sref[II]{II.4.4.7}.

Pongamos
\[
  \sh{G}_1=\Im(\sh{F}\to\sh{G})
  =\Ker(\sh{G}\to\sh{H}),
\]
que es coherente. Como se tiene la sucesión exacta
\[
  0\to\sh{G}_1(n)\to\sh{G}(n)\to\sh{H}(n)
\]
y el funtor $f_*$ es exacto por la izquierda, la sucesión
\[
  0\to f_*(\sh{G}_1(n))\to f_*(\sh{G}(n))
  \to f_*(\sh{H}(n))
\]
es exacta para todo $n$; basta, pues, demostrar que, para $n$ suficientemente grande, la sucesión
\[
  f_*(\sh{F}(n))\to f_*(\sh{G}_1(n))\to0
\]
es exacta. En otras palabras, podemos reducirnos al caso en que $\sh{H}=0$ y el soporte de $\sh{G}$ es propio sobre $Y$, y por tanto \emph{cerrado en $Z$} (II, 5.4.10). En consecuencia, $\sh{G}'=i_*(\sh{G})$ es un $\sh{O}_Z$-módulo coherente tal que $\sh{G}'|X=\sh{G}$. Sabemos por (I, 9.4.3) que existe un $\sh{O}_Z$-módulo coherente $\sh{F}'$ tal que $\sh{F}'|X=\sh{F}$. Además, como $X$ es abierto en $Z$ y $\Supp(\sh{G})\subset X$ es cerrado en $Z$, se define inmediatamente un homomorfismo sobreyectivo $u':\sh{F}'\to\sh{G}'$ cuya restricción a $X$ es el homomorfismo sobreyectivo dado $u:\sh{F}\to\sh{G}$, tomando $u'|U=0$ para todo abierto $U$ de $Z$ que no interseque $\Supp(\sh{G})$, y $u'|U=u|U$ para todo abierto $U\subset X$.

Como al comienzo de la demostración de (2.2.2), podemos reducirnos al caso en que $Y$ es afín; se trata entonces de demostrar que, para $n$ suficientemente grande,
\[
  \Gamma(u):\Gamma(X,\sh{F}(n))
  \to\Gamma(X,\sh{G}(n))
\]
es sobreyectivo. Se tiene el diagrama conmutativo
\[
  \xymatrix@C=35mm@R=12mm{
    \Gamma(Z,\sh{F}'(n)) \ar[r]^{\Gamma(u')} \ar[d]_v &
    \Gamma(Z,\sh{G}'(n)) \ar[d]^w \\
    \Gamma(X,\sh{F}(n)) \ar[r]_{\Gamma(u)} &
    \Gamma(X,\sh{G}(n)).
  }
\]
La definición de $\sh{G}'$ muestra que $w$ es \emph{biyectivo}; por otra parte, en virtud de (2.2.3), $\Gamma(u')$ es \emph{sobreyectivo} para $n$ suficientemente grande, y lo mismo ocurre, por tanto, con $\Gamma(u)$.
\end{proof}

\begin{env}[2.2.5]
\label{III.2.2.5}
\emph{Observaciones.}

(i) La afirmación (i) de \sref{III.2.2.1} sigue siendo cierta si solo se supone que $Y$ es \emph{localmente noetheriano}.
En efecto, la propiedad es manifiestamente local sobre $Y$.
Además, las hipótesis de \sref{III.2.2.1} implican que, para todo subconjunto abierto $U\subset Y$, la restricción de $f$ a $f^{-1}(U)$ es un morfismo proyectivo
\[
  f^{-1}(U)\longrightarrow U
\]
\sref[II]{II.5.5.5}[iii], y $\sh{L}|_{f^{-1}(U)}$ es amplio para este morfismo \sref[II]{II.4.6.4}.

(ii) La afirmación (iii) de \sref{III.2.2.1} sigue siendo válida, como hemos visto, si solo se supone que $X$ es un esquema cuasicompacto, o un preesquema cuyo espacio subyacente es noetheriano, y que $f:X\to Y$ es cuasicompacto \sref[II]{II.4.6.8}.
Conviene observar, sin embargo, que incluso cuando $Y$ es el \emph{espectro de un cuerpo} $K$ y $f$ es \emph{cuasiproyectivo}, la afirmación (ii) de \sref{III.2.2.1} puede dejar de cumplirse.

Por ejemplo, sea
\[
  X'=\Spec(K[T_0,\ldots,T_r])
\]
y sea $X$ la unión de los subconjuntos abiertos afines $D(T_i)$ de $X'$, para $0\leq i\leq r$.
Como la inmersión $X\to X'$ es cuasicompacta, el morfismo estructural $f:X\to Y$ es cuasiafín \sref[II]{II.5.1.10} y, por tanto, $\sh{O}_X$ es muy amplio respecto de $f$ \sref[II]{II.5.1.6}.
Pero
\[
  \Gamma(X,\sh{O}_X)
  =\bigcap_{i=0}^r(K[T_0,\ldots,T_r])_{T_i}
  =K[T_0,\ldots,T_r]
\]
por \sref[I]{I.8.2.1.1}.
Se deduce de \sref{III.1.4.3.1} y \sref{III.1.1.3.5} que
\[
  \HH^r(X,\sh{O}_X^{\otimes n})
  =\HH^r(X,\sh{O}_X)
  \neq0
\]
para todo $n$.

(iii) Bajo las hipótesis de (2.2.4) relativas a $X$, $Y$, $f$ y $\sh{L}$, es inmediato que, si $\sh{H}$ es un $\sh{O}_X$-módulo coherente cuyo soporte es \emph{propio sobre $Y$}, un razonamiento análogo al de (2.2.4) muestra que existe un entero $N$ tal que, para $n\geq N$,
\[
  \RR^i f_*(\sh{H}(n))=0
\]
para todo $i>0$. Se deduce de la sucesión exacta de cohomología que, si $u:\sh{F}\to\sh{G}$ es un homomorfismo de $\sh{O}_X$-módulos coherentes tal que $\Ker(u)$ y $\Coker(u)$ tienen soportes \emph{propios sobre $Y$}, existe un entero $N$ tal que, para $n\geq N$, el homomorfismo correspondiente
\[
  \RR^i f_*(\sh{F}(n))\to\RR^i f_*(\sh{G}(n))
\]
es \emph{biyectivo} para todo $i>0$.
\end{env}


\subsection{Aplicación a los haces graduados de álgebras y de módulos}
\label{subsection:III.2.3}

\begin{theorem}[2.3.1]
\label{III.2.3.1}
Sea $Y$ un preesquema noetheriano, sea $\sh{S}$ una $\sh{O}_Y$-álgebra cuasicoherente de tipo finito, graduada en grados positivos, sea
\[
  X=\Proj(\sh{S}),
\]
y sea $q:X\to Y$ el morfismo estructural.
Sea $\sh{M}$ un $\sh{S}$-módulo graduado cuasicoherente que satisface la condición \textbf{(TF)}.
Existe entonces un entero $N$ tal que, para $n\geq N$, el homomorfismo canónico \sref[II]{II.8.14.5.1}
\[
  \alpha_n:\sh{M}_n
  \longrightarrow q_*\bigl(\shProj_0(\sh{M}(n))\bigr)
  =q_*\bigl((\shProj(\sh{M}))_n\bigr)
\]
\oldpage[III]{103}
es biyectivo.
En otras palabras, el homomorfismo canónico
\[
  \alpha:\sh{M}\longrightarrow\bbGamma_*(\shProj(\sh{M}))
\]
es un \textbf{(TN)}-isomorfismo.
\end{theorem}

\begin{proof}
Podemos restringirnos al caso en que $\sh{M}$ es un $\sh{S}$-módulo de tipo finito \sref[II]{II.3.4.2}.
Como $Y$ es cuasicompacto, existe un entero $d>0$ tal que $\sh{S}^{(d)}$ está engendrada por el $\sh{O}_Y$-módulo cuasicoherente $\sh{S}_d$, que es de tipo finito \sref[II]{II.3.1.10} y, por tanto, coherente porque $Y$ es noetheriano.
Ahora bien, $\sh{M}$ es la suma directa de los $\sh{M}^{(d,k)}$, para $0\leq k<d$, y cada $\sh{M}^{(d,k)}$ es un $\sh{S}^{(d)}$-módulo cuasicoherente de tipo finito, por \sref[II]{II.2.1.6}[iii].
Como la cuestión es local sobre $Y$, basta manifiestamente demostrar que cada homomorfismo canónico
\[
  \alpha:\sh{M}^{(d,k)}
  \longrightarrow
  \bbGamma_*\bigl((\shProj(\sh{M}))^{(d,k)}\bigr)
\]
es un \textbf{(TN)}-isomorfismo.
Teniendo en cuenta \sref[II]{II.8.14.13}, y en particular el diagrama \sref[II]{II.8.14.13.4}, quedamos así reducidos a demostrar el teorema cuando $\sh{S}$ está engendrada por $\sh{S}_1$ y $\sh{S}_1$ es un $\sh{O}_Y$-módulo coherente.

Como $Y$ es noetheriano, el razonamiento utilizado al comienzo de la demostración de \sref{III.2.2.2} muestra que podemos reducirnos al caso
\[
  Y=\Spec(A),\qquad
  \sh{S}=\widetilde{S},\qquad
  \sh{M}=\widetilde{M},
\]
donde $A$ es un anillo noetheriano, $S_1$ es un $A$-módulo de tipo finito y $M$ es un $S$-módulo graduado de tipo finito.
Mostraremos que basta entonces demostrar el teorema para $M=S$.
En efecto, en el caso general existe una sucesión exacta
\[
  L'\longrightarrow L\longrightarrow M\longrightarrow0,
\]
donde $L$ y $L'$ son sumas directas de módulos graduados de la forma $S(m)$.
Si el resultado se cumple para $M=S$, se cumple también para $M=S(m)$ y, por tanto, para $L$ y $L'$.
Consideremos el diagrama conmutativo
\[
  \xymatrix{
    \widetilde{L'_n}\ar[r]\ar[d]_{\alpha_n} &
    \widetilde{L_n}\ar[r]\ar[d]^{\alpha_n} &
    \widetilde{M_n}\ar[r]\ar[d]^{\alpha_n} &
    0\\
    q_*(\widetilde{L'}(n))\ar[r] &
    q_*(\widetilde{L}(n))\ar[r] &
    q_*(\widetilde{M}(n))\ar[r] &
    0 .
  }
\]
\oldpage[III]{104}
La segunda fila es exacta por \sref{III.2.2.3} cuando $n$ es suficientemente grande.
Lo mismo ocurre con la primera fila y, como las dos flechas verticales de la izquierda son isomorfismos, también lo es la tercera.

Queda, pues, demostrar el teorema para $M=S$.
Supongamos primero que
\[
  S=A[T_0,\ldots,T_r],
\]
donde los $T_i$ son indeterminadas.
En este caso, la afirmación es precisamente \sref{III.2.1.11}[ii].
En el caso general, $S$ se identifica con un cociente de un anillo
\[
  S'=A[T_0,\ldots,T_r]
\]
por un ideal graduado, de modo que $X$ es un subesquema cerrado de
\[
  X'=\bb{P}_A^r
\]
\sref[II]{II.2.9.2}.
Si $j:X\to X'$ es la inyección canónica, entonces $j_*(\widetilde{S}(n))$ es el $\sh{O}_{X'}$-módulo
\[
  (\shProj(\widetilde{S}))(n),
\]
donde $S$ se considera como $S'$-módulo graduado; esto se deduce inmediatamente de \sref[II]{II.2.8.7}.
Como $j_*(\widetilde{S}(n))$ es un $\sh{O}_{X'}$-módulo que satisface \textbf{(TF)}, el homomorfismo canónico
\[
  \alpha_n:S_n\longrightarrow
  \Gamma\bigl(X',j_*(\widetilde{S}(n))\bigr)
\]
es biyectivo para todo $n$ suficientemente grande, por el caso polinómico que acabamos de demostrar.
Esto termina la demostración, puesto que
\[
  \Gamma\bigl(X',j_*(\widetilde{S}(n))\bigr)
  =\Gamma(X,\widetilde{S}(n)).
\]
\end{proof}

\begin{corollary}[2.3.2]
\label{III.2.3.2}
Bajo las hipótesis de \sref{III.2.3.1}, sea
\[
  \sh{S}_X=\bigoplus_{n\in\bb{Z}}\sh{O}_X(n),
\]
y sea $\sh{F}$ un $\sh{S}_X$-módulo graduado cuasicoherente de tipo finito.
Entonces $\bbGamma_*(\sh{F})$ satisface la condición \textbf{(TF)}.
\end{corollary}

\begin{proof}
En la demostración de \sref{III.2.3.1} vimos que $X$, que es isomorfo a $\Proj(\sh{S}^{(d)})$ \sref[II]{II.3.1.8}, es de tipo finito sobre $Y$ \sref[II]{II.3.4.1}.
Se deduce de \sref[II]{II.8.14.9} que $\sh{F}$ es isomorfo a un $\sh{S}_X$-módulo graduado de la forma $\shProj(\sh{M})$, donde $\sh{M}$ es un $\sh{S}$-módulo graduado cuasicoherente de tipo finito.
Por \sref{III.2.3.1}, $\bbGamma_*(\sh{F})$ es \textbf{(TN)}-isomorfo a $\sh{M}$ y, por consiguiente, satisface \textbf{(TF)}.
\end{proof}

\begin{env}[2.3.3]
\label{III.2.3.3}
\emph{Escolio.}
Sea $Y$ un preesquema noetheriano, sea $\sh{S}$ una $\sh{O}_Y$-álgebra graduada que satisface las hipótesis de \sref{III.2.3.1}, y sea $X=\Proj(\sh{S})$.
Sea $\cat{K}_{\sh{S}}$ la categoría abeliana de los $\sh{S}$-módulos graduados cuasicoherentes que satisfacen \textbf{(TF)}, y sea $\cat{K}'_{\sh{S}}$ la subcategoría de $\cat{K}_{\sh{S}}$ formada por los $\sh{S}$-módulos que satisfacen \textbf{(TN)}.
Por último, sea $\cat{K}_X$ la categoría de los $\sh{S}_X$-módulos graduados cuasicoherentes $\sh{F}$ de tipo finito; como $\sh{S}_X$ es periódica \sref[II]{II.8.14.4} y \sref[II]{II.8.14.12}, esto equivale a exigir que los $\sh{F}_i$ sean $\sh{O}_X$-módulos coherentes.

Por \sref[II]{II.8.14.8}, \sref[II]{II.8.14.10} y \sref{III.2.3.2}, los funtores
\[
  \sh{M}\longmapsto\shProj(\sh{M}),
  \qquad
  \sh{F}\longmapsto\bbGamma_*(\sh{F})
\]
definen una equivalencia (T, I, 1.2) entre la categoría cociente
\[
  \cat{K}_{\sh{S}}/\cat{K}'_{\sh{S}}
\]
(T, I, 1.11) y la categoría $\cat{K}_X$.
Cuando $\sh{S}$ está engendrada por $\sh{S}_1$, la categoría $\cat{K}_X$ puede sustituirse por la categoría de los $\sh{O}_X$-módulos coherentes \sref[II]{II.8.14.12}.
\end{env}

\begin{proposition}[2.3.4]
\label{III.2.3.4}
Sea $Y$ un preesquema noetheriano.

(i) Sea $\sh{S}$ una $\sh{O}_Y$-álgebra cuasicoherente de tipo finito, graduada en grados positivos.
Pongamos
\[
  X=\Proj(\sh{S}),\qquad
  \sh{S}_X=\shProj(\sh{S})
  =\bigoplus_{n\in\bb{Z}}\sh{O}_X(n).
\]
Entonces $\sh{S}_X$ es una $\sh{O}_X$-álgebra graduada cuasicoherente periódica \sref[II]{II.8.14.12}, sus componentes homogéneas
\[
  (\sh{S}_X)_n=\sh{O}_X(n)
\]
son $\sh{O}_X$-módulos coherentes y, si $d>0$ es un período de $\sh{S}_X$, entonces
\[
  (\sh{S}_X)_d=\sh{O}_X(d)
\]
es $Y$-amplio como $\sh{O}_X$-módulo invertible.
Además, el homomorfismo canónico
\[
  \alpha:\sh{S}\longrightarrow\bbGamma_*(\sh{S}_X)
\]
es un \textbf{(TN)}-isomorfismo.

(ii) Recíprocamente, sea $q:X\to Y$ un morfismo proyectivo y sea $\sh{S}'$ una $\sh{O}_X$-álgebra graduada cuyas componentes homogéneas $\sh{S}'_n$, para $n\in\bb{Z}$, son $\sh{O}_X$-módulos coherentes.
Supongamos que $\sh{S}'$ tiene un período $d>0$ tal que $\sh{S}'_d$ es un $\sh{O}_X$-módulo invertible amplio para $q$.
Entonces
\[
  \sh{S}=\bigoplus_{n\geq0}q_*(\sh{S}'_n)
\]
es una $\sh{O}_Y$-álgebra cuasicoherente de tipo finito, graduada en grados positivos, y existe un $Y$-isomorfismo
\[
  r:X\simto\Proj(\sh{S})
\]
tal que $r^*(\shProj(\sh{S}))$ es isomorfo a $\sh{S}'$ como $\sh{O}_X$-álgebra graduada.
\end{proposition}

\begin{proof}
(i) Casi todas las afirmaciones ya han sido demostradas; la última es un caso particular de \sref{III.2.3.2}.
La periodicidad de $\sh{S}_X$ fue demostrada en \sref[II]{II.8.14.14}, y la afirmación de que existe un período $d>0$ para el cual $\sh{O}_X(d)$ es invertible y $Y$-amplio es precisamente \sref[II]{II.4.6.18}.
Finalmente, para $0\leq k<d$, $(\sh{S}_X)^{(d,k)}$ es un $(\sh{S}_X)^{(d)}$-módulo de tipo finito \sref[II]{II.8.14.14}.
Se deduce que cada $(\sh{S}_X)_n$ es un $\sh{O}_X$-módulo cuasicoherente de tipo finito \sref[II]{II.2.1.6}[ii]; como $\sh{O}_X$ es coherente, también lo es cada $(\sh{S}_X)_n$.

(ii) Tras sustituir el período $d$ por uno de sus múltiplos, podemos suponer que
\[
  \sh{L}=\sh{S}'_d
\]
es muy amplio respecto de $q$ \sref[II]{II.4.6.11}.
Por hipótesis,
\[
  \sh{S}'^{(d)}
  =\bigoplus_{n\in\bb{Z}}\sh{L}^{\otimes n},
\]
y, por tanto,
\[
  \sh{S}^{(d)}
  =\bigoplus_{n\geq0}q_*(\sh{L}^{\otimes n}).
\]
Por \sref[II]{II.3.1.8} y \sref[II]{II.3.2.9}, existe un $Y$-isomorfismo
\[
  s:X'=\Proj(\sh{S})\simto
  X''=\Proj(\sh{S}^{(d)})
\]
tal que
\oldpage[III]{105}
\[
  s^*(\sh{O}_{X''}(n))=\sh{O}_{X'}(nd)
\]
para todo $n$ \sref[II]{II.3.5.2}.
Estableceremos, por tanto, la existencia de un $Y$-isomorfismo $X\simto X'$ demostrando la afirmación siguiente.

\begin{env}[2.3.4.1]
\label{III.2.3.4.1}
Sea $Y$ un preesquema noetheriano, sea $q:X\to Y$ un morfismo proyectivo y sea $\sh{L}$ un $\sh{O}_X$-módulo invertible muy amplio para $q$.
Entonces
\[
  \sh{S}=\bigoplus_{n\geq0}q_*(\sh{L}^{\otimes n})
\]
es una $\sh{O}_Y$-álgebra graduada cuasicoherente de tipo finito tal que
\[
  \sh{S}_n=\sh{S}_1^n
\]
para todo $n$ suficientemente grande, y existe un $Y$-isomorfismo
\[
  r:X\simto P=\Proj(\sh{S})
\]
tal que
\[
  \sh{L}=r^*(\sh{O}_P(1)).
\]
\end{env}

Como $q$ es proyectivo, se deduce de \sref[II]{II.5.4.4} y \sref[II]{II.4.4.7} que existe un $Y$-isomorfismo
\[
  r':X\simto P'=\Proj(\sh{T}),
\]
donde $\sh{T}$ es una $\sh{O}_Y$-álgebra cuasicoherente, $\sh{T}_1$ es de tipo finito y engendra $\sh{T}$, y
\[
  \sh{L}=r'^*(\sh{O}_{P'}(1)).
\]
Si $q':P'\to Y$ es el morfismo estructural, entonces
\[
  \sh{S}=\bigoplus_{n\geq0}q'_*(\sh{O}_{P'}(n)).
\]
Por \sref{III.2.3.1}, el homomorfismo canónico
\[
  \alpha_n:\sh{T}_n\longrightarrow
  \sh{S}_n=q'_*(\sh{O}_{P'}(n))
\]
es biyectivo para todo $n$ suficientemente grande.
Como $\sh{T}_n=\sh{T}_1^n$, se deduce \emph{a fortiori} que $\sh{S}_n=\sh{S}_1^n$ para todo $n$ suficientemente grande.

El homomorfismo canónico $\alpha:\sh{T}\to\sh{S}$ de $\sh{O}_Y$-álgebras graduadas es un \textbf{(TN)}-isomorfismo.
Por consiguiente,
\phantomsection\label{III.2.3.5.1}
\[
  \Phi=\Proj(\alpha):\Proj(\sh{S})\longrightarrow\Proj(\sh{T})
\]
es un isomorfismo \sref[II]{II.3.6.1}.
Además, \sref[II]{II.3.5.2} y el hecho de que los $\sh{T}$-módulos graduados obtenidos de $\sh{S}(n)$ por restricción de escalares son \textbf{(TN)}-isomorfos a $\sh{T}(n)$ dan
\[
  \Phi_*(\sh{O}_P(n))=\sh{O}_{P'}(n)
\]
para todo $n$ \sref[II]{II.3.4.2}.
Solo queda mostrar que $\sh{S}$ es de tipo finito sobre $\sh{O}_Y$.
Los módulos
\[
  \sh{S}_n=q'_*(\sh{O}_{P'}(n))
\]
son coherentes por \sref{III.2.2.1}; y si $\sh{S}_n=\sh{S}_1^n$ para $n\geq n_0$, entonces $\sh{S}$ está engendrada por el módulo coherente
\[
  \bigoplus_{i\leq n_0}\sh{S}_i.
\]
Esto demuestra la afirmación \sref[I]{I.9.6.2}.

Volvamos a la demostración de \sref{III.2.3.4}, conservando su notación.
Hemos demostrado que existe un $Y$-isomorfismo
\[
  r'':X\simto X''
\]
tal que
\[
  r''_*(\sh{L}^{\otimes n})=\sh{O}_{X''}(n)
\]
para todo $n\in\bb{Z}$.
Sea $q'':X''\to Y$ el morfismo estructural.
El álgebra $\sh{S}'$ es la suma directa de los $\sh{S}'^{(d)}$-módulos graduados $\sh{S}'^{(d,k)}$.
Cada uno de ellos es un $\sh{S}'^{(d)}$-módulo cuasicoherente de tipo finito, por la periodicidad de $\sh{S}'$ y la hipótesis de que los $\sh{S}'_n$ son $\sh{O}_X$-módulos de tipo finito \sref[II]{II.8.14.12}.
Pongamos
\[
  \sh{F}^{(k)}=r''_*(\sh{S}'^{(d,k)}),
\]
de modo que los $\sh{F}^{(k)}$ son $\sh{S}_{X''}$-módulos graduados cuasicoherentes de tipo finito.
Por \sref[II]{II.8.14.8}, el homomorfismo canónico
\[
  \beta:\shProj\bigl(\bbGamma_*(\sh{F}^{(k)})\bigr)
  \longrightarrow\sh{F}^{(k)}
\]
es un isomorfismo de $\sh{S}_{X''}$-módulos.
Por otra parte,
\[
  q''_*\bigl((\sh{F}^{(k)})_n\bigr)
  =q_*\bigl((\sh{S}'^{(d,k)})_n\bigr),
\]
y para $n\geq0$ el último $\sh{O}_Y$-módulo es, por definición, igual a $(\sh{S}^{(d,k)})_n$.
En otras palabras, la inyección canónica
\[
  \sh{S}^{(d,k)}
  \longrightarrow\bbGamma_*(\sh{F}^{(k)})
\]
es un \textbf{(TN)}-isomorfismo.
Por tanto,
\[
  \shProj(\sh{S}^{(d,k)})
  =\shProj\bigl(\bbGamma_*(\sh{F}^{(k)})\bigr),
\]
y, en consecuencia,
\[
  r''{}^*\bigl(\shProj(\sh{S}^{(d,k)})\bigr)
  \simeq\sh{S}'^{(d,k)}.
\]
Queda observar que, salvo isomorfismo canónico,
\[
  \shProj(\sh{S}^{(d,k)})
  =s_*\bigl((\shProj(\sh{S}))^{(d,k)}\bigr)
\]
\sref[II]{II.8.14.13.1}; esto demuestra que $r^*(\shProj(\sh{S}))$ es isomorfo a $\sh{S}'$.

Finalmente, por \sref{III.2.3.2}, cada $\bbGamma_*(\sh{F}^{(k)})$ satisface la condición \textbf{(TF)} y, por tanto, también la satisface cada $\sh{S}^{(d,k)}$.
Además, los $\sh{S}_n=q_*(\sh{S}'_n)$ son coherentes por \sref{III.2.2.1}, porque los $\sh{S}'_n$ son coherentes.
Se deduce inmediatamente que los $\sh{S}^{(d,k)}$ son $\sh{S}^{(d)}$-módulos de tipo finito.
Como \sref{III.2.3.4.1} muestra que $\sh{S}^{(d)}$ es una $\sh{O}_Y$-álgebra de tipo finito, $\sh{S}$ es asimismo una $\sh{O}_Y$-álgebra de tipo finito.
\end{proof}

\oldpage[III]{106}
\begin{proposition}[2.3.5]
\label{III.2.3.5}
Sea $Y$ un preesquema noetheriano integral, sea $X$ un preesquema integral y sea
\[
  f:X\longrightarrow Y
\]
un morfismo proyectivo birracional.
Existe entonces un ideal fraccionario coherente
\[
  \sh{J}\subset\sh{R}(Y)
\]
\sref[II]{II.8.1.2} tal que $X$ es $Y$-isomorfo al preesquema obtenido por explosión de $\sh{J}$ \sref[II]{II.8.1.3}.
Además, existe un subconjunto abierto $U$ de $Y$ tal que la restricción de $f$ a $f^{-1}(U)$ es un isomorfismo de $f^{-1}(U)$ sobre $U$ (cf.\ \textbf{I}, 6.5.5), y tal que $\sh{J}|U$ es invertible.
\end{proposition}

\begin{proof}
Existe un $\sh{O}_X$-módulo invertible $\sh{L}$ muy amplio para $f$ \sref[II]{II.4.4.2} y \sref[II]{II.5.3.2}.
Aplicando \sref{III.2.3.4.1}, identificamos $X$ con $\Proj(\sh{S})$, donde
\[
  \sh{S}=\bigoplus_{n\geq0}f_*(\sh{L}^{\otimes n}).
\]
Los módulos $f_*(\sh{L}^{\otimes n})$ son sin torsión sobre $\sh{O}_Y$ \sref[I]{I.7.4.5}.
Por tanto, $\sh{S}$ se identifica canónicamente con un sub-$\sh{O}_Y$-módulo de
\[
  \sh{S}\otimes_{\sh{O}_Y}\sh{R}(Y)
\]
\sref[I]{I.7.4.1}.
Este último es un haz simple \sref[I]{I.7.3.6} y está, por tanto, determinado por su restricción a un subconjunto abierto no vacío.

Elijamos un subconjunto abierto no vacío $U'\subset U$ sobre el cual $\sh{L}|f^{-1}(U')$ sea isomorfo a $\sh{O}_X|f^{-1}(U')$ y, por tanto, a $\sh{O}_Y|U'$.
Las restricciones $f_*(\sh{L}^{\otimes n})|U'$ son entonces isomorfas a $\sh{O}_Y|U'$.
Se sigue que
\[
  \sh{S}\otimes_{\sh{O}_Y}\sh{R}(Y)
  \simeq\sh{R}(Y)[T],
\]
donde $T$ es una indeterminada.
Por \sref{III.2.3.4.1}, $\sh{S}$ es \textbf{(TN)}-isomorfa a la sub-$\sh{O}_Y$-álgebra engendrada por la imagen canónica de $f_*(\sh{L})$ en esta álgebra de polinomios.
Bajo la identificación mostrada, dicha imagen es $\sh{J}\cdot T$, donde $\sh{J}$ es un sub-$\sh{O}_Y$-módulo coherente de $\sh{R}(Y)$ \sref{III.2.2.1}.
Su restricción a $U'$ es isomorfa a $\sh{O}_Y|U'$ y, reduciendo $U$ si es necesario, $\sh{J}|U$ es invertible.
Por consiguiente,
\[
  \sh{S}\quad\text{es \textbf{(TN)}-isomorfa a}\quad
  \bigoplus_{n\geq0}\sh{J}^n,
\]
lo que termina la demostración.
\end{proof}

\begin{corollary}[2.3.6]
\label{III.2.3.6}
Bajo las hipótesis de \sref{III.2.3.5}, supongamos además que para todo sub-$\sh{O}_Y$-módulo coherente no nulo
\[
  \sh{J}\subset\sh{R}(Y)
\]
existe un $\sh{O}_Y$-módulo invertible $\sh{L}$ tal que
\[
  \Gamma\bigl(Y,\sh{L}\otimes_{\sh{O}_Y}
  \shHom(\sh{J},\sh{O}_Y)\bigr)\neq0.
\]
Entonces, en el enunciado de \sref{III.2.3.5}, puede tomarse $\sh{J}$ como un ideal de $\sh{O}_Y$.
Esta condición adicional se satisface siempre si existe un $\sh{O}_Y$-módulo amplio.
\end{corollary}

\begin{proof}
En efecto, por \sref[0]{0.5.4.2},
\[
  \sh{L}\otimes\shHom(\sh{J},\sh{O}_Y)
  =\shHom\bigl(\sh{L}^{-1},\shHom(\sh{J},\sh{O}_Y)\bigr)
  =\shHom(\sh{J}\otimes\sh{L}^{-1},\sh{O}_Y).
\]
La hipótesis dice, por tanto, que existe un homomorfismo no nulo
\[
  u:\sh{J}\otimes\sh{L}^{-1}\longrightarrow\sh{O}_Y.
\]
Para todo $y\in Y$, el módulo $(\sh{J}\otimes\sh{L}^{-1})_y$ se identifica con un sub-$\sh{O}_y$-módulo del cuerpo de fracciones $\sh{R}(Y)_y$ de $\sh{O}_y$ \sref[I]{I.7.1.5}.
Así, $u_y$ es necesariamente inyectivo, y $u$ identifica $\sh{J}\otimes\sh{L}^{-1}$ con un ideal $\sh{J}'$ de $\sh{O}_Y$.
Pero
\[
  \Proj\left(\bigoplus_{n\geq0}\sh{J}^n\right)
  \quad\text{y}\quad
  \Proj\left(\bigoplus_{n\geq0}
  (\sh{J}\otimes\sh{L}^{-1})^n\right)
\]
son $Y$-isomorfos \sref[II]{II.3.1.8}; esto demuestra la primera aserción.

Para la segunda, observemos que
\[
  \sh{F}=\shHom_{\sh{O}_Y}(\sh{J},\sh{O}_Y)
\]
es coherente y no nulo, pues existe un subconjunto abierto $U$ de $Y$ sobre el cual $\sh{J}|U$ es invertible.
Si $\sh{L}$ es un $\sh{O}_Y$-módulo amplio, existe un entero $n$ tal que
\[
  \sh{F}(n)=\sh{F}\otimes\sh{L}^{\otimes n}
\]
está engendrado por sus secciones sobre $Y$ \sref[II]{II.4.5.5}.
\emph{A fortiori},
\[
  \Gamma(Y,\sh{F}(n))\neq0,
\]
lo que demuestra la aserción.
\end{proof}

\begin{corollary}[2.3.7]
\label{III.2.3.7}
Sean $X$ e $Y$ esquemas integrales proyectivos sobre un cuerpo $k$, y sea
\[
  f:X\longrightarrow Y
\]
un $k$-morfismo birracional.
Entonces $X$ es $k$-isomorfo a un $Y$-esquema obtenido por explosión de un subesquema cerrado $Y'$ de $Y$, no necesariamente reducido.
\end{corollary}

\begin{proof}
\oldpage[III]{107}
En efecto, $f$ es proyectivo \sref[II]{II.5.5.5}[v], y como $Y$ es proyectivo sobre $k$, se satisface la condición adicional de \sref{III.2.3.6}.
Basta, por tanto, tomar el subesquema cerrado $Y'$ de $Y$ definido por el ideal coherente $\sh{J}$ de \sref{III.2.3.6}.
\end{proof}

\begin{env}[2.3.8]
\label{III.2.3.8}
\emph{Observación.}
En el Capítulo~IV, al estudiar la noción de divisor, veremos que si, en el enunciado de \sref{III.2.3.5}, se supone que los anillos locales $\sh{O}_y$, para $y\in Y$, son factoriales---como ocurre, por ejemplo, cuando $Y$ es no singular---, entonces $X$ puede obtenerse de $Y$ por explosión de un subpreesquema cerrado $Y'$ de $Y$ cuyo espacio subyacente está contenido en $Y-U$.
\end{env}


\subsection{Una generalización del teorema fundamental}
\label{subsection:III.2.4}

\begin{theorem}[2.4.1]
\label{III.2.4.1}
Sea $Y$ un preesquema noetheriano, sea $\sh{S}$ una $\sh{O}_Y$-álgebra cuasicoherente de tipo finito, sea
\[
  f:X\longrightarrow Y
\]
un morfismo proyectivo, pongamos $\sh{S}'=f^*(\sh{S})$, y sea $\sh{M}$ un $\sh{S}'$-módulo cuasicoherente de tipo finito.
Entonces:
\begin{enumerate}
  \item[(i)] Para todo $p\in\bb Z$, $\RR^p f_*(\sh{M})$ es un $\sh{S}$-módulo de tipo finito.

  \item[(ii)] Supongamos además que $\sh{L}$ es un $\sh{O}_X$-módulo invertible amplio para $f$ y, para todo $n\in\bb Z$, pongamos
  \[
    \sh{M}(n)=\sh{M}\otimes\sh{L}^{\otimes n}.
  \]
  Existe un entero $N$ tal que, para $n\geq N$,
  \[
  \label{III.2.4.1.1}
    \RR^p f_*(\sh{M}(n))=0
    \tag{2.4.1.1}
  \]
  para todo $p>0$, y el homomorfismo canónico
  \[
    f^*\bigl(f_*(\sh{M}(n))\bigr)\longrightarrow\sh{M}(n)
  \]
  \sref[0]{0.4.4.3} es sobreyectivo.
\end{enumerate}
\end{theorem}

\begin{proof}
Pongamos
\[
  Y'=\Spec(\sh{S}),\qquad X'=\Spec(\sh{S}'),
\]
de modo que $X'=X\times_Y Y'$ \sref[II]{II.1.5.5}.
Sean
\[
  g:Y'\longrightarrow Y,\qquad g':X'\longrightarrow X
\]
los morfismos estructurales, que son afines por definición, y sea
\[
  f'=f_{(Y')}:X'\longrightarrow Y'.
\]
Tenemos así un diagrama conmutativo
\[
  \xymatrix{
    X\ar[d]_f &
    X'\ar[l]_{g'}\ar[d]^{f'}\ar[dl]^{h}\\
    Y &
    Y'\ar[l]^{g},
  }
\]
donde $f'$ es proyectivo \sref[II]{II.5.5.5}[iii] y
\[
  h=f\circ g'=g\circ f'.
\]

\emph{(i)} Sea $\widetilde{\sh{M}}$ el $\sh{O}_{X'}$-módulo asociado al $\sh{S}'$-módulo cuasicoherente $\sh{M}$ cuando se considera $X'$ como un $X$-esquema afín \sref[II]{II.1.4.3}; así,
\[
  \sh{M}=g'_*(\widetilde{\sh{M}}).
\]
Como $\sh{M}$ es un $\sh{S}'$-módulo de tipo finito, $\widetilde{\sh{M}}$ es un $\sh{O}_{X'}$-módulo de tipo finito \sref[II]{II.1.4.5}.
Además, $h$ es de tipo finito, puesto que $g$ y $f'$ son de tipo finito \sref[II]{II.1.3.7} y \sref[I]{I.6.3.4}[ii].
Por tanto, $X'$ es noetheriano \sref[I]{I.6.3.7}, y $\widetilde{\sh{M}}$ es, en consecuencia, coherente.

Como $g'$ es afín, el homomorfismo canónico
\[
  \RR^p f_*(\sh{M})
  \longrightarrow
  \RR^p h_*(\widetilde{\sh{M}})
\]
es biyectivo \sref{III.1.3.4}.
Es, además, un homomorfismo de $\sh{S}$-módulos.
En efecto, el homomorfismo canónico
\[
\label{III.2.4.1.2}
  g'_*(\sh{O}_{X'})
  \otimes_{\sh{O}_X}
  g'_*(\widetilde{\sh{M}})
  \longrightarrow
  g'_*(\widetilde{\sh{M}})
  \tag{2.4.1.2}
\]
define la estructura de $\sh{S}'$-módulo sobre $\sh{M}$, puesto que
\[
  \sh{S}'=g'_*(\sh{O}_{X'}).
\]
Da lugar canónicamente a un homomorfismo
\[
  f_*\bigl(g'_*(\sh{O}_{X'})\bigr)
  \otimes
  \RR^p f_*\bigl(g'_*(\widetilde{\sh{M}})\bigr)
  \longrightarrow
  \RR^p f_*\bigl(g'_*(\widetilde{\sh{M}})\bigr).
\]
\oldpage[III]{108}
Por \sref[0]{0.12.2.2}, y dado que \eqref{III.2.4.1.2} se obtiene a su vez, aplicando \sref[0]{0.4.2.2.1}, del homomorfismo
\[
  \sh{O}_{X'}\otimes\widetilde{\sh{M}}
  \longrightarrow\widetilde{\sh{M}}
\]
que define la estructura de $\sh{O}_{X'}$-módulo sobre $\widetilde{\sh{M}}$, el diagrama
\[
  \xymatrix{
    f_*\bigl(g'_*(\sh{O}_{X'})\bigr)
    \otimes\RR^p f_*\bigl(g'_*(\widetilde{\sh{M}})\bigr)
    \ar[r]\ar[d] &
    \RR^p f_*\bigl(g'_*(\widetilde{\sh{M}})\bigr)\ar[d]\\
    h_*(\sh{O}_{X'})
    \otimes\RR^p h_*(\widetilde{\sh{M}})
    \ar[r] &
    \RR^p h_*(\widetilde{\sh{M}})
  }
\]
es conmutativo \sref[0]{0.12.2.6}.
Componiendo las flechas horizontales con el homomorfismo inducido por el homomorfismo canónico
\[
  \sh{S}
  \longrightarrow f_*(f^*(\sh{S}))
  =f_*(\sh{S}')
  =f_*\bigl(g'_*(\sh{O}_{X'})\bigr)
  =h_*(\sh{O}_{X'}),
\]
obtenemos la aserción.

Por otra parte, como $g$ es afín y $f'$ es separado y cuasicompacto, el homomorfismo canónico
\[
  \RR^p h_*(\widetilde{\sh{M}})
  \longrightarrow
  g_*\bigl(\RR^p f'_*(\widetilde{\sh{M}})\bigr)
\]
es biyectivo \sref{III.1.4.14}.
Como antes, es un isomorfismo de $\sh{S}$-módulos, esta vez utilizando la conmutatividad de \sref[0]{0.12.2.6.2}.
Ahora bien, $f'$ es proyectivo y $\widetilde{\sh{M}}$ es coherente, por lo que $\RR^p f'_*(\widetilde{\sh{M}})$ es un $\sh{O}_{Y'}$-módulo coherente por \sref{III.2.2.1}.
Se sigue que
\[
  g_*\bigl(\RR^p f'_*(\widetilde{\sh{M}})\bigr)
\]
es un $\sh{S}$-módulo de tipo finito \sref[II]{II.1.4.5}.

\emph{(ii)} Pongamos
\[
  \sh{L}'=g'^*(\sh{L}),
\]
un $\sh{O}_{X'}$-módulo invertible.
Para todo $n\in\bb Z$, salvo isomorfismo canónico,
\[
  g'_*\bigl(\widetilde{\sh{M}}\otimes\sh{L}'^{\otimes n}\bigr)
  =
  g'_*(\widetilde{\sh{M}})\otimes\sh{L}^{\otimes n}
  =
  \sh{M}(n)
\]
\sref[0]{0.5.4.10}.
Aplicando el razonamiento de \emph{(i)} a
\[
  \widetilde{\sh{M}}\otimes\sh{L}'^{\otimes n}
\]
muestra que
\[
  \RR^p f_*
  \bigl(g'_*(\widetilde{\sh{M}}\otimes\sh{L}'^{\otimes n})\bigr)
  \simeq
  g_*\bigl(\RR^p f'_*
  (\widetilde{\sh{M}}\otimes\sh{L}'^{\otimes n})\bigr).
\]
Ahora bien, $\sh{L}'$ es amplio para $f'$ \sref[II]{II.4.6.13}[iii].
Se deduce de \sref{III.2.2.1} que existe un entero $N$ tal que
\[
  \RR^p f'_*
  \bigl(\widetilde{\sh{M}}\otimes\sh{L}'^{\otimes n}\bigr)=0
\]
para todo $p>0$ y todo $n\geq N$.
Esto demuestra \eqref{III.2.4.1.1}.

Se deduce una vez más de \sref{III.2.2.1} que puede elegirse $N$ de modo que, para $n\geq N$, el homomorfismo canónico
\[
  f'^*\bigl(f'_*(\widetilde{\sh{M}}\otimes\sh{L}'^{\otimes n})\bigr)
  \longrightarrow
  \widetilde{\sh{M}}\otimes\sh{L}'^{\otimes n}
\]
sea sobreyectivo.
Como $g'_*$ es exacto \sref[II]{II.1.4.4}, el homomorfismo correspondiente
\[
  g'_*\Bigl(
    f'^*\bigl(f'_*(\widetilde{\sh{M}}\otimes\sh{L}'^{\otimes n})\bigr)
  \Bigr)
  \longrightarrow
  g'_*(\widetilde{\sh{M}}\otimes\sh{L}'^{\otimes n})
  =\sh{M}(n)
\]
es sobreyectivo.
Pero \sref[II]{II.1.5.2} da
\[
  g'_*\Bigl(
    f'^*\bigl(f'_*(\widetilde{\sh{M}}\otimes\sh{L}'^{\otimes n})\bigr)
  \Bigr)
  =
  f^*\Bigl(
    g_*\bigl(f'_*(\widetilde{\sh{M}}\otimes\sh{L}'^{\otimes n})\bigr)
  \Bigr),
\]
y $g_*\circ f'_*=f_*\circ g'_*$.
Por consiguiente,
\[
  g'_*\Bigl(
    f'^*\bigl(f'_*(\widetilde{\sh{M}}\otimes\sh{L}'^{\otimes n})\bigr)
  \Bigr)
  =
  f^*\bigl(f_*(\sh{M}(n))\bigr),
\]
lo que termina la demostración.
\end{proof}

\begin{env}[2.4.2]
\label{III.2.4.2}
Aplicaremos \sref{III.2.4.1}, en particular, cuando $\sh{S}$ sea una $\sh{O}_Y$-álgebra graduada en grados positivos y
\[
  \sh{M}=\bigoplus_{k\in\bb Z}\sh{M}_k
\]
sea un $\sh{S}'$-módulo graduado.
\oldpage[III]{109}
Bajo las mismas hipótesis de finitud sobre $\sh{S}$ y $\sh{M}$, se deduce de \sref{III.2.4.1} que
\[
  \bigoplus_{k\in\bb Z}\RR^p f_*(\sh{M}_k)
\]
es un $\sh{S}$-módulo de tipo finito para todo $p$.
Bajo las hipótesis adicionales de \sref{III.2.4.1}[ii], existe además un entero $N$ tal que, para $n\geq N$,
\[
  \RR^p f_*(\sh{M}_k(n))=0
\]
para todo $p>0$ y todo $k\in\bb Z$, y el homomorfismo canónico
\[
  f^*\bigl(f_*(\sh{M}_k(n))\bigr)
  \longrightarrow
  \sh{M}_k(n)
\]
es sobreyectivo para todo $k\in\bb Z$.
\end{env}


\subsection{Característica de Euler--Poincar\'e y polinomio de Hilbert}
\label{subsection:III.2.5}

\begin{env}[2.5.1]
\label{III.2.5.1}
Sea $A$ un anillo artiniano y sea $X$ un $A$-esquema proyectivo sobre
\[
  Y=\Spec(A).
\]
Para todo $\sh{O}_X$-módulo coherente $\sh{F}$, los $A$-módulos
\[
  \HH^i(X,\sh{F})\qquad(i\geq0)
\]
son de tipo finito \sref{III.2.2.1} y, por tanto, tienen longitud finita, puesto que $A$ es artiniano.
Además, \sref{III.2.2.1} muestra que $\HH^i(X,\sh{F})=0$ salvo para un número finito de $i\geq0$.
Así, el entero
\[
\label{III.2.5.1.1}
  \chi_A(\sh{F})
  =
  \sum_{i=0}^{\infty}
  (-1)^i\operatorname{length}_A\bigl(\HH^i(X,\sh{F})\bigr)
  \tag{2.5.1.1}
\]
está definido para todo $\sh{O}_X$-módulo coherente $\sh{F}$.
Cuando $A$ es un anillo local artiniano, $\chi_A(\sh{F})$ se llama la \emph{característica de Euler--Poincar\'e} de $\sh{F}$ con respecto a $A$.
Para $\sh{F}=\sh{O}_X$, el entero $\chi_A(\sh{O}_X)$ se llama el \emph{género aritmético} de $X$ con respecto a $A$.
\end{env}

\begin{proposition}[2.5.2]
\label{III.2.5.2}
Sea
\[
  0\longrightarrow\sh{F}'
  \longrightarrow\sh{F}
  \longrightarrow\sh{F}''
  \longrightarrow0
\]
una sucesión exacta de $\sh{O}_X$-módulos coherentes.
Entonces
\[
\label{III.2.5.2.1}
  \chi_A(\sh{F})
  =
  \chi_A(\sh{F}')+\chi_A(\sh{F}'').
  \tag{2.5.2.1}
\]
\end{proposition}

\begin{proof}
Los módulos de cohomología de $\sh{F}$, $\sh{F}'$ y $\sh{F}''$ se anulan en todos los grados suficientemente grandes.
Existe, por tanto, un entero $r>0$ para el cual la sucesión exacta larga de cohomología tiene la forma
\[
  0\longrightarrow\HH^0(X,\sh{F}')
  \longrightarrow\HH^0(X,\sh{F})
  \longrightarrow\HH^0(X,\sh{F}'')
  \longrightarrow\HH^1(X,\sh{F}')
  \longrightarrow\cdots
\]
\[
  \cdots\longrightarrow\HH^r(X,\sh{F}')
  \longrightarrow\HH^r(X,\sh{F})
  \longrightarrow\HH^r(X,\sh{F}'')
  \longrightarrow0.
\]
En una sucesión exacta de $A$-módulos de longitud finita con términos nulos en ambos extremos, la suma alternada de las longitudes es nula \sref[0]{0.11.10.1}.
La aplicación de este hecho da inmediatamente \eqref{III.2.5.2.1}.
\end{proof}

El resultado de \sref{III.2.5.2} se aplica siempre que $X$ sea un $A$-esquema cuasicompacto y que los $A$-módulos $\HH^i(X,\sh{F})$ sean de tipo finito para todo $\sh{O}_X$-módulo coherente $\sh{F}$ \sref{III.1.4.12}.

\begin{theorem}[2.5.3]
\label{III.2.5.3}
Sea $A$ un anillo local artiniano, sea $X$ un esquema proyectivo sobre
\[
  Y=\Spec(A),
\]
sea $\sh{L}$ un $\sh{O}_X$-módulo invertible muy amplio relativo a $Y$, y sea $\sh{F}$ un $\sh{O}_X$-módulo coherente.
Para todo $n\in\bb Z$, pongamos
\[
  \sh{F}(n)=\sh{F}\otimes_{\sh{O}_X}\sh{L}^{\otimes n}.
\]
\begin{enumerate}
  \item[(i)] Existe un único polinomio $P\in\bb Q[T]$ tal que
  \[
    \chi_A(\sh{F}(n))=P(n)
  \]
  para todo $n\in\bb Z$.
  El polinomio $P$ se llama el \emph{polinomio de Hilbert} de $\sh{F}$ con respecto a $A$.

  \item[(ii)] Para todo $n$ suficientemente grande,
  \[
    \chi_A(\sh{F}(n))
    =
    \operatorname{length}_A\Gamma(X,\sh{F}(n)).
  \]

  \item[(iii)] El coeficiente del término de mayor grado de $\chi_A(\sh{F}(n))$ es no negativo.
\end{enumerate}
\end{theorem}

Además, demostraremos en el Capítulo~IV, en la sección dedicada a la dimensión, que \emph{el grado de $\chi_A(\sh{F}(n))$ es igual a la dimensión del soporte de $\sh{F}$}.

\begin{proof}
Para todo $n$ suficientemente grande, se tiene
\[
  \HH^i(X,\sh{F}(n))=0
\]
para todo $i>0$ \sref{III.2.2.1}.
\oldpage[III]{110}
Por consiguiente,
\[
  \chi_A(\sh{F}(n))
  =
  \operatorname{length}_A\HH^0(X,\sh{F}(n))
  =
  \operatorname{length}_A\Gamma(X,\sh{F}(n))
\]
para todo $n$ suficientemente grande, lo que demuestra \emph{(ii)}.
Se sigue que $\chi_A(\sh{F}(n))\geq0$ para todo $n$ suficientemente grande y, por tanto, \emph{(iii)} se deduce de \emph{(i)}.
La aserción de unicidad de \emph{(i)} es inmediata, por lo que queda demostrar la existencia de $P$.

Mostremos primero que podemos suponer
\[
  \mathfrak{m}\sh{F}=0,
\]
donde $\mathfrak{m}$ es el ideal maximal de $A$.
En efecto, existe un entero $s>0$ tal que $\mathfrak{m}^s=0$ y, por tanto, $\sh{F}(n)$ posee la filtración finita
\[
  \sh{F}(n)
  \supset\mathfrak{m}\sh{F}(n)
  \supset\cdots
  \supset\mathfrak{m}^{s-1}\sh{F}(n)
  \supset0.
\]
Una inducción que utiliza \eqref{III.2.5.2.1} da
\[
  \chi_A(\sh{F}(n))
  =
  \sum_{k=1}^{s}
  \chi_A\left(
    \mathfrak{m}^{k-1}\sh{F}(n)/
    \mathfrak{m}^{k}\sh{F}(n)
  \right).
\]
Pero
\[
  \mathfrak{m}^{k-1}\sh{F}(n)/
  \mathfrak{m}^{k}\sh{F}(n)
  =
  \sh{F}'_k(n),
  \qquad
  \sh{F}'_k=
  \mathfrak{m}^{k-1}\sh{F}/\mathfrak{m}^{k}\sh{F},
\]
lo que demuestra la reducción.

Supongamos, pues, que $\mathfrak{m}\sh{F}=0$.
Sea $X'$ el subesquema cerrado de $X$ obtenido como imagen inversa, por el morfismo estructural
\[
  X\longrightarrow\Spec(A),
\]
del único punto cerrado de $\Spec(A)$, y sea $j:X'\to X$ la inyección canónica.
Entonces
\[
  \sh{F}=j_*(\sh{F}')
\]
para un $\sh{O}_{X'}$-módulo coherente $\sh{F}'$, y $X'$ es proyectivo sobre $\Spec(K)$, donde
\[
  K=A/\mathfrak{m}.
\]
Si $\sh{L}'=j^*(\sh{L})$, entonces $\sh{L}'$ es muy amplio relativo a $\Spec(K)$ \sref[II]{II.4.4.10}, y
\[
  \sh{F}(n)=j_*(\sh{F}'(n)),
  \qquad
  \sh{F}'(n)
  =
  \sh{F}'\otimes_{\sh{O}_{X'}}\sh{L}'^{\otimes n}
\]
\sref[0]{0.5.4.10}.
Se sigue que
\[
  \chi_A(\sh{F}(n))
  =
  \chi_K(\sh{F}'(n))
\]
(G, II, 4.9.1), y quedamos reducidos al caso en que $A$ es un cuerpo.

Ahora puede considerarse $X$ como un subesquema cerrado de
\[
  P=\mathbf{P}_A^r
\]
para un $r$ conveniente \sref[II]{II.5.5.4}[ii].
Si $i:X\to P$ es la inyección canónica, el mismo razonamiento da
\[
  \chi_A(\sh{F}(n))
  =
  \chi_A\bigl(i_*(\sh{F})(n)\bigr).
\]
Podemos, por tanto, limitarnos al caso
\[
  X=\mathbf{P}_A^r=\Proj(S),
  \qquad
  S=A[T_0,\ldots,T_r],
\]
donde $A$ es un cuerpo.

En este caso, $\sh{F}=\widetilde{M}$ para un $S$-módulo graduado $M$ de tipo finito \sref[II]{II.2.7.8}.
Por el teorema de las sicigias de Hilbert (M, VIII, 6.5), $M$ posee una resolución finita por $S$-módulos libres graduados de tipo finito:
\[
  0\longrightarrow L_q
  \longrightarrow L_{q-1}
  \longrightarrow\cdots
  \longrightarrow L_1
  \longrightarrow M
  \longrightarrow0.
\]
Como el funtor $M\mapsto\widetilde{M}$ es exacto \sref[II]{II.2.5.4}, obtenemos una sucesión exacta
\[
  0\longrightarrow\widetilde{L}_q
  \longrightarrow\widetilde{L}_{q-1}
  \longrightarrow\cdots
  \longrightarrow\widetilde{L}_1
  \longrightarrow\widetilde{M}
  \longrightarrow0,
\]
y, por tanto, para todo $n\in\bb Z$, una sucesión exacta
\[
  0\longrightarrow\widetilde{L}_q(n)
  \longrightarrow\widetilde{L}_{q-1}(n)
  \longrightarrow\cdots
  \longrightarrow\widetilde{L}_1(n)
  \longrightarrow\widetilde{M}(n)
  \longrightarrow0.
\]
Una inducción sobre $q$, utilizando \sref{III.2.5.2}, da
\[
  \chi_A(\widetilde{M}(n))
  =
  \sum_{j=1}^{q}
  (-1)^{j+1}\chi_A(\widetilde{L}_j(n)).
\]
Para demostrar \emph{(i)}, quedamos, por tanto, reducidos al caso en que $M$ es graduado, libre y de tipo finito y, por consiguiente, al caso
\[
  M=S(h)
\]
para algún $h\in\bb Z$.
Entonces
\[
  \widetilde{M}(n)
  =
  \widetilde{M(n)}
  =
  \widetilde{S(n+h)}
\]
\sref[II]{II.2.5.15}, de modo que el teorema se deduce del lema siguiente.
\end{proof}

\oldpage[III]{111}
\begin{lemma}[2.5.3.1]
\label{III.2.5.3.1}
Sea $A$ un cuerpo, sea $r>0$ un entero y sea
\[
  X=\mathbf{P}_A^r.
\]
Entonces
\[
  \chi_A(\sh{O}_X(n))
  =
  \binom{n+r}{r}
\]
para todo $n\in\bb Z$.
\end{lemma}

\begin{proof}
Para $n\geq0$,
\[
  \chi_A(\sh{O}_X(n))
  =
  \operatorname{length}_A\HH^0(X,\sh{O}_X(n)),
\]
que es el número de monomios en los $T_i$ de grado total $n$, a saber,
\[
  \binom{n+r}{r}
\]
\sref{III.2.1.12}.

Para $n\leq-r-1$, se tiene análogamente
\[
  \chi_A(\sh{O}_X(n))
  =
  (-1)^r\operatorname{length}_A\HH^r(X,\sh{O}_X(n)).
\]
Escribiendo $n=-r-h$, la dimensión de $\HH^r(X,\sh{O}_X(n))$ sobre $A$ es el número de sucesiones
\[
  (p_i)_{0\leq i\leq r}
\]
de enteros positivos que satisfacen
\[
  \sum_{i=0}^{r}p_i=r+h
\]
\sref{III.2.1.12}.
Equivalentemente, es el número de sucesiones de enteros no negativos
\[
  (q_i)_{0\leq i\leq r}
\]
tales que
\[
  \sum_{i=0}^{r}q_i=h-1.
\]
Este número es
\[
  \binom{h+r-1}{r}
  =
  (-1)^r\binom{n+r}{r}.
\]
Finalmente, si $-r\leq n<0$, entonces
\[
  \binom{n+r}{r}=0,
\]
y $\HH^i(X,\sh{O}_X(n))=0$ para todo $i\geq0$ \sref{III.2.1.12}.
Esto demuestra el lema.
\end{proof}

\begin{corollary}[2.5.4]
\label{III.2.5.4}
Sea $A$ un anillo local artiniano, sea $S$ una $A$-álgebra graduada de tipo finito engendrada por $S_1$, sea $M$ un $S$-módulo graduado de tipo finito y sea
\[
  X=\Proj(S).
\]
Entonces, para todo $n$ suficientemente grande,
\[
  \chi_A(\widetilde{M}(n))
  =
  \operatorname{length}_A(M_n).
\]
\end{corollary}

\begin{proof}
Para todo $n$ suficientemente grande, los módulos
\[
  M_n
  \quad\text{y}\quad
  \Gamma(X,\widetilde{M}(n))
\]
son isomorfos \sref{III.2.3.1}.
\end{proof}


\subsection{Aplicación: criterios de amplitud}
\label{subsection:III.2.6}

\begin{proposition}[2.6.1]
\label{III.2.6.1}
Sea $Y$ un preesquema noetheriano, sea $f:X\to Y$ un morfismo propio,
y sea $\sh{L}$ un $\sh{O}_X$-módulo invertible.
Las condiciones siguientes son equivalentes:
\begin{enumerate}
  \item[\emph{a)}] $\sh{L}$ es amplio para $f$.
  \item[\emph{b)}] Para todo $\sh{O}_X$-módulo coherente $\sh{F}$, existe
  un entero $N$ tal que, para $n\geq N$,
  \[
    \RR^qf_*(\sh{F}\otimes\sh{L}^{\otimes n})=0
  \]
  para todo $q>0$.
  \item[\emph{c)}] Para todo ideal coherente $\sh{J}$ de $\sh{O}_X$, existe
  un entero $N$ tal que, para $n\geq N$,
  \[
    \RR^1f_*(\sh{J}\otimes\sh{L}^{\otimes n})=0.
  \]
\end{enumerate}
\end{proposition}

\begin{proof}
Hemos visto que \emph{a)} implica \emph{b)} \sref{III.2.2.1},~(ii).
Es inmediato que \emph{b)} implica \emph{c)}, y queda demostrar
que \emph{c)} implica \emph{a)}.
Podemos limitarnos al caso en que $Y$ es afín \sref[II]{II.4.6.4}.
Basta mostrar que, cuando $h$ recorre las secciones de
$\sh{L}^{\otimes n}$ sobre $X$, con $n>0$, los subconjuntos abiertos afines $X_h$
forman un recubrimiento de $X$ \sref[II]{II.4.5.2}.

Para ello, sea $x$ un punto cerrado de $X$ y sea $U$ un entorno
abierto afín de $x$.
Encontraremos un entero $n$ y una sección
\[
  h\in\Gamma(X,\sh{L}^{\otimes n})
\]
tal que $x\in X_h\subset U$.
El subconjunto abierto $X_h$ será entonces afín \sref[I]{I.1.3.6}.
La unión de todos esos $X_h$ es un subconjunto abierto de $X$ que contiene todo
punto cerrado y, por tanto, es todo $X$, puesto que $X$ es noetheriano
\sref[I]{I.6.3.7} y \sref[0]{0.2.1.3}.

Sean $\sh{J}$ y $\sh{J}'$, respectivamente, los ideales cuasicoherentes de
$\sh{O}_X$ que definen los subpreesquemas cerrados reducidos cuyos espacios
subyacentes son $X-U$ y $(X-U)\cup\{x\}$, respectivamente \sref[I]{I.5.2.1}.
Los ideales $\sh{J}$ y $\sh{J}'$ son coherentes \sref[I]{I.6.1.1},
$\sh{J}'\subset\sh{J}$, y
\[
  \sh{J}''=\sh{J}/\sh{J}'
\]
es un $\sh{O}_X$-módulo coherente \sref[0]{0.5.3.3}, con soporte
$\{x\}$, y $\sh{J}''_x=\kres(x)$.
Como $\sh{L}$ es localmente libre, la sucesión
\[
  0\longrightarrow
  \sh{J}'\otimes\sh{L}^{\otimes n}
  \longrightarrow
  \sh{J}\otimes\sh{L}^{\otimes n}
  \longrightarrow
  \sh{J}''\otimes\sh{L}^{\otimes n}
  \longrightarrow0
\]
es exacta para todo $n$.
Por hipótesis, para todo $n$ suficientemente grande,
\[
  \HH^1(X,\sh{J}'\otimes\sh{L}^{\otimes n})=0.
\]

\oldpage[III]{112}

La sucesión exacta de cohomología muestra, por tanto, que
\[
  \Gamma(X,\sh{J}\otimes\sh{L}^{\otimes n})
  \longrightarrow
  \Gamma(X,\sh{J}''\otimes\sh{L}^{\otimes n})
\]
es sobreyectivo.
Una sección $g$ de $\sh{J}''\otimes\sh{L}^{\otimes n}$ sobre $X$ tal que
$g(x)\neq0$ es, por consiguiente, la imagen de una sección
\[
  h\in\Gamma(X,\sh{J}\otimes\sh{L}^{\otimes n})
  \subset\Gamma(X,\sh{L}^{\otimes n});
\]
la inclusión se deduce de \sref[0]{0.5.4.1}.
Por construcción, $h(x)\neq0$, mientras que $h(z)=0$ para todo $z\notin U$.
Así, $x\in X_h\subset U$, lo que demuestra la aserción.
\end{proof}

\begin{proposition}[2.6.2]
\label{III.2.6.2}
Sea $Y$ un preesquema noetheriano, sea $f:X\to Y$ un morfismo de
tipo finito, sea $g:X'\to X$ un morfismo finito sobreyectivo, sea
$\sh{L}$ un $\sh{O}_X$-módulo invertible, y pongamos
\[
  \sh{L}'=g^*(\sh{L}).
\]
Supongamos que existe un subconjunto $Z$ de $X$, propio sobre $Y$
\sref[II]{II.5.4.10}, tal que, para todo $x\in X-Z$, o bien $X$ es
normal en $x$, o bien
\[
  \bigl(g_*(\sh{O}_{X'})\bigr)_x
\]
es un $\sh{O}_{X,x}$-módulo libre.
Entonces $\sh{L}$ es amplio para $f$ si y solo si $\sh{L}'$ es amplio para
$f\circ g$.
\end{proposition}

\begin{proof}
\begin{env}[2.6.2.1]
\label{III.2.6.2.1}
Como $g$ es afín, la condición es necesaria \sref[II]{II.5.1.12}.
Para demostrar que es suficiente, podemos suponer que $Y$ es afín
\sref[II]{II.4.6.4}.
Mostremos primero que también podemos limitarnos al caso en que $X$ es
reducido.

Sea $j:X_{\mathrm{red}}\to X$ la inmersión canónica, y pongamos
\[
  X_1=X_{\mathrm{red}},
  \qquad
  X'_1=X'\times_X X_1.
\]
Tenemos entonces el diagrama conmutativo
\[
  \xymatrix{
    X'\ar[d]_g &
    X'_1\ar[l]_{j'}\ar[d]^{g_1}\\
    X &
    X_1\ar[l]^{j}
  }
  \label{III.2.6.2.2}\tag{2.6.2.2}
\]
El morfismo $f\circ j$ es de tipo finito \sref[I]{I.6.3.4}, y
$g_1$ es finito \sref[II]{II.6.1.5},~(iii).
Si $\sh{L}'$ es amplio para $f\circ g$, entonces
$j'^*(\sh{L}')$ es amplio para $f\circ g\circ j'$, porque $j'$ es una
inmersión cerrada \sref[II]{II.5.1.12} y \sref[I]{I.4.3.2}.
Si $Z_1=j^{-1}(Z)$, entonces $Z_1$ es propio sobre $Y$
\sref[II]{II.5.4.10}.
Además, si $X$ es normal en un punto $x$, también lo es $X_{\mathrm{red}}$
en ese punto; y si
$\bigl(g_*(\sh{O}_{X'})\bigr)_x$ es un $\sh{O}_{X,x}$-módulo libre,
entonces \sref[II]{II.1.5.2} muestra que
\[
  \bigl((g_1)_*(\sh{O}_{X'_1})\bigr)_x
\]
es un $\sh{O}_{X_1,x}$-módulo libre.
Finalmente, como $X$ es noetheriano \sref[I]{I.6.3.7}, la amplitud de
$j^*(\sh{L})$ implica la de $\sh{L}$ \sref[II]{II.4.5.14}.
Como
\[
  j'^*(\sh{L}')=g_1^*(j^*(\sh{L})),
\]
esto demuestra la reducción anunciada.
En lo sucesivo suponemos que $Y$ es afín y que $X$ es reducido.
\end{env}

Se satisfacen ahora las hipótesis de \sref[II]{II.6.6.11}.
Existen un $Y$-preesquema reducido $X_2$ y un
$Y$-morfismo finito birracional
\[
  h:X_2\to X
\]
cuya restricción a $h^{-1}(X-Z)$ es un isomorfismo sobre $X-Z$, y
tal que $h^*(\sh{L})$ es amplio.
Reemplazando $X'$ por $X_2$, quedamos reducidos a demostrar la proposición bajo
la hipótesis adicional de que $g$ posee las propiedades que acabamos de enumerar para
$h$.
Seguimos denotando por $Z$ un subpreesquema cerrado de $X$, propio sobre
$Y$ \sref[II]{II.5.4.10}, cuyo espacio subyacente es el subconjunto dado
$Z$.

\begin{env}[2.6.2.3]
\label{III.2.6.2.3}
Sea $X_1$ un subpreesquema cerrado de $X$, sea $j:X_1\to X$ la
inmersión canónica, y sea
\[
  X'_1=g^{-1}(X_1)=X'\times_X X_1,
\]
con inmersión canónica $j':X'_1\to X'$.
El cuadrado correspondiente es de nuevo \eqref{III.2.6.2.2}.
Pongamos
\[
  \sh{L}_1=j^*(\sh{L}),
  \qquad
  \sh{L}'_1=j'^*(\sh{L}')=g_1^*(\sh{L}_1).
\]
Entonces $\sh{L}'_1$ es amplio para $f\circ g\circ j'$
\sref[II]{II.5.1.12}.
Si $Z_1=j^{-1}(Z)$, el subpreesquema cerrado $Z_1$ de $X_1$ es propio
sobre $Y$ \sref[II]{II.5.4.2},~(ii).
En otras palabras, las hipótesis de \sref{III.2.6.2} se cumplen para
$X_1$, $\sh{L}_1$, $g_1$ y $Z_1$.
\end{env}

\oldpage[III]{113}

Esto nos permite demostrar \sref{III.2.6.2} por inducción noetheriana
\sref[0]{0.2.2.2} cuando la restricción de $g$ a $g^{-1}(X-Z)$ es un
isomorfismo sobre $X-Z$; como se vio en \sref{III.2.6.2.1}, ese caso es
suficiente.
Basta mostrar que, si la conclusión de \sref{III.2.6.2} se cumple
para $\sh{L}_1$ sobre todo subpreesquema cerrado $X_1$ de $X$ cuyo
espacio subyacente no es $X$, entonces también se cumple para $\sh{L}$.

\begin{env}[2.6.2.4]
\label{III.2.6.2.4}
Pongamos
\[
  \sh{A}=\sh{O}_X,
  \qquad
  \sh{B}=g_*(\sh{O}_{X'}).
\]
Como $g$ es birracional y $X$ es reducido, $\sh{B}$ es una
sub-$\sh{A}$-álgebra del haz $\sh{R}(X)$ de anillos totales de fracciones,
y es un $\sh{A}$-módulo coherente.
Además,
\[
  \sh{B}|_{X-Z}=\sh{A}|_{X-Z}.
\]
Sea $\sh{K}$ el conductor de $\sh{B}$ sobre $\sh{A}$, es decir, el mayor
sub-$\sh{A}$-módulo cuasicoherente de $\sh{A}$ tal que
\[
  \sh{B}\sh{K}\subset\sh{A}.
\]
Equivalentemente, $\sh{K}$ es el anulador del
$\sh{A}$-módulo $\sh{B}/\sh{A}$ \sref[0]{0.5.3.7}.
Se sigue que
\[
  \sh{B}\sh{K}=\sh{K}.
\]
Está claro que $\sh{K}_x=\sh{A}_x$ en todo punto que admite un entorno
$W_x$ sobre el cual $g:g^{-1}(W_x)\to W_x$ es un isomorfismo.
Esto ocurre, en particular, en todo punto de $X-Z$, y en un
entorno de todo punto genérico de una componente irreducible de
$X$.

Sea
\[
  Z_1=\Spec(\sh{A}/\sh{K})
\]
el subpreesquema cerrado de $X$ definido por $\sh{K}$.
Es propio sobre $Y$, porque su espacio subyacente es cerrado en $Z$
\sref[II]{II.5.4.10}.
La definición de $\sh{K}$ da también
\[
  \sh{B}|_{X-Z_1}=\sh{A}|_{X-Z_1}.
\]
Podemos, por tanto, reducirnos al caso
\[
  Z=\Spec(\sh{A}/\sh{K}).
\]
Como $X-Z_1$ es un subconjunto abierto no vacío de $X$, podemos además
suponer que el espacio subyacente de $Z$ es distinto de $X$.
\end{env}

\begin{env}[2.6.2.5]
\label{III.2.6.2.5}
Consideremos $X'$ como $\Spec(\sh{B})$, puesto que $g$ es afín, y sea
\[
  \sh{K}'=\widetilde{\sh{K}}.
\]
Este es un ideal coherente de $\sh{O}_{X'}$ que satisface
$g_*(\sh{K}')=\sh{K}$ \sref[II]{II.1.4.1}.
El subpreesquema cerrado
\[
  Z'=g^{-1}(Z)=Z\times_X X'
\]
de $X'$ está definido por $\sh{K}'$ y es igual a
\[
  \Spec(\sh{B}/\sh{K})
\]
\sref[II]{II.1.4.10}.
El morfismo $h:Z'\to Z$ es finito \sref[II]{II.6.1.5},~(iii);
por consiguiente, $Z'$ es propio sobre $Y$ \sref[II]{II.6.1.11} y
\sref[II]{II.5.4.2},~(ii).

Debemos demostrar que, para todo $x\in X$ y todo entorno abierto $U$
de $x$, existen un entero $n>0$ y una sección $s$ de
$\sh{L}^{\otimes n}$ sobre $X$ tales que
\[
  x\in X_s\subset U
\]
\sref[II]{II.4.5.2}.
Distinguimos dos casos.

\smallskip
\noindent
$1^\circ$ Supongamos primero que $x\in X-Z$.
Podemos suponer evidentemente que $U\subset X-Z$, de modo que
$U'=g^{-1}(U)$ no corta a $Z'$.
Como $\sh{L}'$ es amplio, existen $n>0$ y una sección $s'$ de
$\sh{L}'^{\otimes n}$ sobre $X'$ tales que
\[
  x'=g^{-1}(x)\in X'_{s'}\subset g^{-1}(U)
\]
\sref[II]{II.4.5.2}.
También podemos suponer que
$\sh{K}'\otimes\sh{L}'^{\otimes n}$ está engendrado por sus secciones globales
\sref[II]{II.4.5.5}.
Como $\sh{K}'_{x'}=\sh{O}_{X',x'}$, una de estas secciones, digamos $s''$, es
no nula en $x'$.
Multiplicándola por $s'$ y reemplazando así $n$ por $2n$, podemos conseguir
que
\[
  x'\in X'_{s''}\subset g^{-1}(U).
\]
Por \sref[0]{0.5.4.10}, existe un isomorfismo canónico
\[
  \Gamma(X,\sh{K}\otimes\sh{L}^{\otimes n})
  \xrightarrow{\sim}
  \Gamma(X',\sh{K}'\otimes\sh{L}'^{\otimes n}).
\]
La sección $s$ de $\sh{K}\otimes\sh{L}^{\otimes n}$ correspondiente a
$s''$ posee manifiestamente las propiedades requeridas.

\smallskip
\noindent
$2^\circ$ Supongamos ahora que $x\in Z$.
Sea $\sh{J}$ el ideal coherente de $\sh{O}_X$ que define el
subpreesquema cerrado reducido cuyo espacio subyacente es $X-U$, y consideremos en
$\sh{B}$ los ideales coherentes

\oldpage[III]{114}

\[
  \sh{J}\sh{B}
  \qquad\text{y}\qquad
  \sh{J}\sh{K}
  =
  \sh{J}(\sh{K}\sh{B})
  =
  \sh{K}(\sh{J}\sh{B}).
\]
Tenemos el diagrama de inclusiones
\[
  \xymatrix@R=1.5em{
    \sh{J}\sh{B}\ar@{^{(}->}[r] &
    \sh{B}\\
    \sh{J}\ar@{^{(}->}[r]\ar@{^{(}->}[u] &
    \sh{A}\ar@{^{(}->}[u]\\
    \sh{J}\sh{K}\sh{B}=\sh{J}\sh{K}
      \ar@{^{(}->}[r]\ar@{^{(}->}[u] &
    \sh{K}\ar@{^{(}->}[u]
  }
  \label{III.2.6.2.6}\tag{2.6.2.6}
\]
Sea $\sh{J}'=(\sh{J}\sh{B})^\sim$ el ideal coherente de
$\sh{O}_{X'}$ tal que
\[
  g_*(\sh{J}')=\sh{J}\sh{B}.
\]
Entonces
\[
  \sh{J}'\sh{K}'=(\sh{J}\sh{K}\sh{B})^\sim,
  \qquad
  \sh{J}'/\sh{J}'\sh{K}'
  =
  (\sh{J}\sh{B}/\sh{J}\sh{K}\sh{B})^\sim
\]
\sref[II]{II.1.4.4}.
Como $\sh{J}|_V=\sh{J}\sh{K}|_V$ sobre todo subconjunto abierto $V$ disjunto
de $Z$, el soporte de $\sh{J}'/\sh{J}'\sh{K}'$ está contenido en
$Z'$.
Como $Z'$ es propio sobre $Y$, \sref{III.2.2.4} muestra que, para todo
$n$ suficientemente grande, la aplicación canónica
\[
  \Gamma(X',\sh{J}'\otimes\sh{L}'^{\otimes n})
  \longrightarrow
  \Gamma\bigl(
    X',
    (\sh{J}'/\sh{J}'\sh{K}')\otimes\sh{L}'^{\otimes n}
  \bigr)
\]
es sobreyectiva.
Por \sref[0]{0.5.4.10}, se sigue que la aplicación canónica
\[
  \Gamma(X,\sh{J}\sh{B}\otimes\sh{L}^{\otimes n})
  \longrightarrow
  \Gamma\bigl(
    X,
    (\sh{J}\sh{B}/\sh{J}\sh{K}\sh{B})
      \otimes\sh{L}^{\otimes n}
  \bigr)
\]
es sobreyectiva.

Sean $i:Z\to X$ e $i':Z'\to X'$ las inmersiones canónicas.
Forman parte del diagrama conmutativo
\[
  \xymatrix{
    X'\ar[d]_g &
    Z'\ar[l]_{i'}\ar[d]^h\\
    X &
    Z\ar[l]^{i}
  }
\]
Pongamos
\[
  \sh{M}=i^*(\sh{L}),
  \qquad
  \sh{M}'=i'^*(\sh{L}').
\]
Como $\sh{L}'$ es amplio, $\sh{M}'$ es amplio
\sref[II]{II.5.1.12}, y
\[
  \sh{M}'=h^*(\sh{M}).
\]
La hipótesis de inducción noetheriana, aplicada puesto que $Z\neq X$, muestra ahora
que $\sh{M}$ es amplio.
Por consiguiente, para todo $n$ suficientemente grande,
\[
  i^*(\sh{J}/\sh{J}\sh{K})\otimes\sh{M}^{\otimes n}
\]
está engendrado por sus secciones sobre $Z$ \sref[II]{II.4.5.5}.
Como
\[
  \sh{J}/\sh{J}\sh{K}
  =
  i_*\bigl(i^*(\sh{J}/\sh{J}\sh{K})\bigr),
\]
\sref[0]{0.5.4.10} da una sección $s$ de
\[
  (\sh{J}/\sh{J}\sh{K})\otimes\sh{L}^{\otimes n}
\]
sobre $X$, para algún $n$ suficientemente grande, tal que $s(x)\neq0$.
En efecto, $\sh{J}_x=\sh{A}_x$ por la definición de $\sh{J}$, mientras que
$\sh{K}_x\neq\sh{A}_x$ por hipótesis.

El diagrama \eqref{III.2.6.2.6} muestra que $s$ es también una sección de
\[
  (\sh{J}\sh{B}/\sh{J}\sh{K}\sh{B})
  \otimes\sh{L}^{\otimes n}
\]
sobre $X$.
Es, por tanto, la imagen canónica de una sección
\[
  t\in\Gamma(X,\sh{J}\sh{B}\otimes\sh{L}^{\otimes n}).
\]
Por definición, la imagen de $t$ módulo
$\sh{J}\sh{K}\sh{B}\otimes\sh{L}^{\otimes n}$ pertenece a
\[
  \frac{\sh{J}\otimes\sh{L}^{\otimes n}}
       {\sh{J}\sh{K}\sh{B}\otimes\sh{L}^{\otimes n}}.
\]
El diagrama \eqref{III.2.6.2.6} implica, por tanto, que $t$ es una sección
de $\sh{J}\otimes\sh{L}^{\otimes n}$ y, en particular, una sección de
$\sh{L}^{\otimes n}$.
Se tiene $t(x)\neq0$, luego $x\in X_t$.
Por la definición de $\sh{J}$, la sección $t$ se anula sobre $X-U$, el
soporte de $\sh{O}_X/\sh{J}$.
Así, $X_t\subset U$, y la demostración está terminada.
\end{env}
\end{proof}

\oldpage[III]{115}

\begin{env}[2.6.3]
\label{III.2.6.3}
\emph{Observación.}
Cuando $X$ es propio sobre $Y$, la Proposición~\sref{III.2.6.2} puede demostrarse
más sencillamente razonando como en el teorema de Chevalley
\sref[II]{II.6.7.1}, con ayuda de la Proposición~\sref{III.2.6.1} y
del lema \sref[II]{II.6.7.1.1}.
\end{env}
